\input{preamble} 

% OK, start here.
%
\begin{document} 

\title{Cycles relatifs} 


\maketitle

\phantomsection
\label{section-phantom} 

\tableofcontents



\section{Introduction} 
\label{section-introduction} 

\noindent
 Une référence fondamentale est \cite{SV}. 

\medskip\noindent
 Dans ce chapitre, nous ne définissons que les cycles relatifs universellement
 entiers de \cite{SV}. Ce choix permet de développer une théorie 
un peu plus simple que dans l'article original, au prix toutefois d'une certaine perte. 

\medskip\noindent
 Fixons un morphisme $X \to S$ de type fini 
entre schémas noethériens. Une famille $\alpha$ de $r$-cycles sur les fibres
 de $X/S$ est simplement une famille $\alpha = (\alpha_s)_{s \in S}$
 , où $\alpha_s \in Z_r(X_s)$. Le changement 
de base $g^*\alpha$ de $\alpha$ se définit immédiatement pour tout morphisme $g : S' \to S$.
 On dit alors que $\alpha$ est un {\it $r$-cycle relatif sur $X/S$}
 si $\alpha$ est compatible avec les spécialisations : pour tout 
morphisme $g : S' \to S$, où $S'$ est le spectre d'un anneau 
de valuation discrète, la fibre générique de $g^*\alpha$
 doit se spécialiser en la fibre fermée de $g^*\alpha$.
 Voir la section \ref{section-families-specialization}. 






\section{Conventions et notation} 
\label{section-conventions} 

\noindent
 On consultera le chapitre Homologie de Chow et classes de Chern pour
 nos conventions et notations relatives aux cycles sur les schémas 
localement de type fini sur une base noethérienne fixée ; voir 
Homologie de Chow, section \ref{chow-section-setup} et suivantes. 

\medskip\noindent
 En particulier, si  $X$  est localement de type fini sur un corps $k$, 
alors  $Z_r(X)$  désigne le groupe des cycles de dimension $r$, voir
 Homologie de Chow, exemple  \ref{chow-example-field}  et 
section \ref{chow-section-cycles}. Étant donné un sous-schéma fermé
 intègre  $Z \subset X$  avec $\dim(Z) = r$ , nous avons  $[Z] \in Z_r(X)$
  et si  $X$  est quasi-compact, alors $Z_r(X)$  est abélien libre sur ces classes. 




\section{Cycles relatifs aux corps} 
\label{section-relative-fields} 

\noindent
 Soit $k$ un corps. Soit  $X$  un schéma localement algébrique sur  $k$. 
Soit  $r \geq 0$  un entier. Dans ce cadre, nous avons le groupe 
$Z_r(X)$  des  $r$-cycles sur  $X$, voir la section  \ref{section-conventions}. 

\medskip\noindent
 {\bf  Changement de base. }  Pour toute extension de corps  $k'/k$  il existe une application de changement 
de base  $Z_r(X) \to Z_r(X_{k'})$, voir
 Homologie de Chow, section  \ref{chow-section-change-base}. À 
savoir, étant donné un sous-schéma fermé intègre  $Z \subset X$
  de dimension  $r$  on envoie  $[Z] \in Z_r(X)$  au $r$-cycle  
$[Z_{k'}]_r \in Z_r(X_{k'})$  associé au sous-schéma fermé  
$Z_{k'} \subset X_{k'}$  (bien sûr, en général  $Z_{k'}$
  n’est ni irréductible ni réduit). L’application de changement de base  
$Z_r(X) \to Z_r(X_{k'})$  est toujours injective. 

\begin{lemma} 
\label{lemma-multiplicities-field-extension} 
Soit $K/k$ une extension de corps. Soit $Z$ un schéma intègre localement 
algébrique sur $k$. La multiplicité  $m_{Z', Z_K}$  d’une composante 
irréductible  $Z' \subset Z_K$  est  $1$  ou une puissance de la caractéristique de  $k$. 
\end{lemma} 

\begin{proof}
 Si la caractéristique de  $k$  est zéro, alors  $k$  est parfait et 
la multiplicité est toujours  $1$  puisque $X_K$ est réduit d'après
 Variétés, lemme  \ref{varieties-lemma-geometrically-reduced}. 
Supposons que la caractéristique de  $k$ soit  $p > 0$. 
Soit  $L$  le corps des fonctions de  $Z$. Puisque  $Z$  est localement algébrique
 sur  $k$, l'extension de corps $L/k$ est finiment engendrée. 
L'anneau  $K \otimes_k L$  est noethérien
 (Algèbre, lemme  \ref{algebra-lemma-Noetherian-field-extension}).
 En termes algébriques, il faut montrer que la longueur de l'anneau
 local artinien  $(K \otimes_k L)_\mathfrak q$
  est une puissance de  $p$ pour chaque idéal premier minimal  $\mathfrak q$. 

\medskip\noindent
 Soit  $L'/L$  une extension finie purement inséparable, disons de degré 
$p^n$. Alors  $K \otimes_k L \subset K \otimes_k L'$  est un homomorphisme
 d'anneaux fini et libre de rang $p^n$  qui induit un homéomorphisme sur les 
spectres et des extensions de corps résiduels purement inséparables. 
Ainsi, pour chaque idéal premier minimal $\mathfrak q$  comme ci-dessus,
 il existe un unique idéal premier minimal 
$\mathfrak q' \subset K \otimes_k L'$  au-dessus de celui-ci et 
$$
p^n \text{length}((K \otimes_k L)_\mathfrak q) =
[\kappa(\mathfrak q') : \kappa(\mathfrak q)]
\text{length}((K \otimes_k L')_{\mathfrak q'})
$$
 par Algèbre, lemme  \ref{algebra-lemma-pushdown-module}  appliqué 
à  $M = (K \otimes_k L')_{\mathfrak q'} \cong
(K \otimes_k L)_{\mathfrak q}^{\oplus p^n}$. 
Puisque  $[\kappa(\mathfrak q') : \kappa(\mathfrak q)]$  est une puissance
 de  $p$ , nous concluons qu'il suffit de prouver 
l'énoncé pour  $L'$  et  $\mathfrak q'$. 

\medskip\noindent
 D'après le paragraphe précédent et Algèbre, lemme  \ref{algebra-lemma-make-separable} , nous 
pouvons supposer qu'il existe une tour de corps $L/k'/k$ telle que $L/k'$ soit séparable et
 que $k'/k$ soit finie purement inséparable. Alors  $K \otimes_k k'$ est 
un anneau local artinien. L'argument du paragraphe précédent 
(appliqué à $L = k$  et  $L' = k'$) montre que  $\text{length}(K \otimes_k k')$
  est une puissance de $p$. Comme $L/k'$ est la localisation 
d'une $k'$-algèbre 
lisse (Algèbre, lemme \ref{algebra-lemma-localization-smooth-separable}),
 $S = (K \otimes_k L)_\mathfrak q$ est la localisation d'une 
$R = K \otimes_k k'$-algèbre lisse en un idéal premier minimal.
 Donc  $R \to S$  est un homomorphisme local plat d'anneaux 
locaux artiniens et  $\mathfrak m_R S = \mathfrak m_S$. Il
 découle de Algèbre, lemme  \ref{algebra-lemma-pullback-module}  que  
$\text{length}(K \otimes_k k') = \text{length}(R) =
\text{length}(S) = \text{length}((K \otimes_k L)_\mathfrak q)$
  et la preuve est terminée. 
\end{proof} 

\begin{lemma} 
\label{lemma-how-different} 
Soit  $k$  un corps de caractéristique  $p > 0$  avec clôture parfaite $k^{perf}$. 
Soit  $X$  un schéma algébrique sur  $k$. Soit $r \geq 0$  un entier. Le 
conoyau de l'application injective  $Z_r(X) \to Z_r(X_{k^{perf}})$  est un module de torsion de 
$p$-puissance (Compléments sur l'algèbre, Définition 
\ref{more-algebra-definition-f-power-torsion}). 
\end{lemma} 

\begin{proof} 
Puisque  $X$  est quasi-compact, le groupe abélien  $Z_r(X)$  est libre avec une 
base donnée par les sous-schémas fermés intègres de dimension $r$. De même pour  
$Z_r(X_{k^{perf}})$. 
Puisque  $X_{k^{perf}} \to X$  est un homéomorphisme, il en résulte
 que  $Z_r(X) \to Z_r(X_{k^{perf}})$  est injectif avec conoyau de torsion. Chaque
 élément du conoyau est de torsion de $p$-puissance d'après le
 lemme  \ref{lemma-multiplicities-field-extension}. 
\end{proof} 





\section{Spécialisation des cycles} 
\label{section-specialization} 

\noindent
 Soit  $R$  un anneau de valuation discrète avec corps des fractions $K$
  et corps résiduel  $\kappa$. Soit  $X$  un schéma localement de type fini 
sur  $R$. Soit  $r \geq 0$. Il existe une application de spécialisation 
$$
sp_{X/R} : Z_r(X_K) \longrightarrow Z_r(X_\kappa)
$$
 définie comme suit. Pour un sous-schéma fermé et intègre  $Z \subset X_K$
  de dimension $r$ , nous notons $\overline{Z}$ l'image schématique de
  $Z \to X$. On définit alors $sp_{X/R}$ comme l'unique application $\mathbf{Z}$-linéaire
 telle que 
$$
sp_{X/R}([Z]) = [\overline{Z}_\kappa]_r
$$
 Expliquons brièvement pourquoi cette définition est correcte. Tout d'abord, observons que le
 morphisme  $X_K \to X$  est quasi-compact et donc le morphisme  $Z \to X$
 est quasi-compact. La formation de l'image schématique de $Z \to X$
 commute donc aux changements de base plats 
d'après Morphismes, lemme  \ref{morphisms-lemma-flat-base-change-scheme-theoretic-image}. En 
particulier, après changement de base à $X_K$, on obtient $Z = \overline{Z}_K$. 
Puisque  $Z$  est intègre, bien sûr  $\overline{Z}$  est également intègre
 et en fait est égal au sous-schéma fermé intègre unique dont le point
 générique est l'image du point générique de  $Z$. Il résulte de 
Variétés, lemme  \ref{varieties-lemma-dominate-valuation-ring-dimension-fibres}  
que  $Z_\kappa$  est équidimensionnel de dimension $r$. 

\begin{lemma} 
\label{lemma-specialization-module} 
Soit  $R$  un anneau de valuation discrète avec corps des fractions $K$  et corps résiduel  
$\kappa$. Soit  $X$  un schéma localement de type fini sur  $R$. Soit  $r \geq 0$. 
Soit  $\mathcal{F}$  un $\mathcal{O}_X$-module cohérent plat sur  $R$. Supposons  
$\dim(\text{Supp}(\mathcal{F}_K)) \leq r$. Alors 
$\dim(\text{Supp}(\mathcal{F}_\kappa)) \leq r$  et 
$$
sp_{X/R}([\mathcal{F}_K]_r) = [\mathcal{F}_\kappa]_r
$$
\end{lemma} 

\begin{proof} 
L'énoncé sur la dimension découle de Compléments sur les morphismes, lemme 
\ref{more-morphisms-lemma-relative-dimension-support-flat}. 
Soit  $x$  un point générique d'un sous-schéma fermé intègre  
$Z \subset X_\kappa$ de dimension  $r$. Pour terminer la preuve,
 nous montrerons que le coefficient de $[Z]$
 dans le membre de gauche (L) et le membre de droite (R) de l’égalité est le même. 

\medskip\noindent
 Posons $A = \mathcal{O}_{X, x}$ et $M = \mathcal{F}_x$.
 Remarquons que  $M$  est un $A$-module fini plat sur  $R$. 
Soit  $\pi \in R$  un uniformisant de sorte que  
$A/\pi A = \mathcal{O}_{X_\kappa, x}$.
 D'après Homologie de Chow, lemme \ref{chow-lemma-additivity-divisors-restricted} , nous
 avons 
$$
\sum\nolimits_i \text{length}_A(A/(\pi, \mathfrak q_i))
\text{length}_{A_{\mathfrak q_i}}(M_{\mathfrak q_i}) =
\text{length}_A(M/\pi M)
$$
 où la somme porte sur les idéaux premiers minimaux  
$\mathfrak q_i$  dans le support de $M$. 
Puisque  $\pi$  est un non-diviseur de zéro sur  $M$ , nous voyons
 que  $\pi \not \in \mathfrak q_i$  et donc ces
 idéaux premiers correspondent aux points génériques  $y_i \in X_K$  du 
support de  $\mathcal{F}_K$  qui se spécialisent en notre $x \in X_\kappa$ 
choisi. Ainsi, le membre de gauche est le coefficient de  $[Z]$
  dans (L). Bien sûr,  $\text{length}_A(M/\pi M)$  est le coefficient 
de  $[Z]$  dans (R). Cela termine la démonstration. 
\end{proof} 

\begin{lemma} 
\label{lemma-specialization-closed} 
Soit  $R$  un anneau de valuation discrète avec corps des fractions $K$  et corps résiduel  
$\kappa$. Soit  $X$  un schéma localement de type fini sur  $R$. Soit  $r \geq 0$. 
Soit  $W \subset X$  un sous-schéma fermé plat sur  $R$. Supposons  
$\dim(W_K) \leq r$. Alors $\dim(W_\kappa) \leq r$  et 
$$
sp_{X/R}([W_K]_r) = [W_\kappa]_r
$$
\end{lemma} 

\begin{proof} En
 prenant  $\mathcal{F} = \mathcal{O}_W$ , il s'agit d'un cas particulier du
 lemme  \ref{lemma-specialization-module}. Voir 
Homologie de Chow, lemme \ref{chow-lemma-cycle-closed-coherent}. 
\end{proof} 

\begin{lemma} 
\label{lemma-specialization-extension} 
Soit  $R'/R$  une extension d'anneaux de valuation discrète induisant une extension de corps des
 fractions  $K'/K$  et une extension de corps résiduels $\kappa'/\kappa$
  (Compléments sur l'algèbre, définition  
\ref{more-algebra-definition-extension-discrete-valuation-rings}). 
Soit  $X$ localement de type fini sur  $R$. On note  $X' = X_{R'}$.
 Alors le diagramme 
$$
\xymatrix{
Z_r(X'_{K'}) \ar[rr]_{sp_{X'/R'}} & & Z_r(X'_{\kappa'}) \\
Z_r(X_K) \ar[rr]^{sp_{X/R}} \ar[u] & & Z_r(X_\kappa) \ar[u]
}
$$
 est commutatif, où  $r \geq 0$  et les flèches verticales sont des applications de changement de base. 
\end{lemma} 

\begin{proof} 
On observe que  $X'_{K'} = X_{K'} = X_K \times_{\Spec(K)} \Spec(K')$
  et de même pour les fibres fermées, de sorte que les flèches verticales ont effectivement
 un sens (voir la section  \ref{section-relative-fields}).
 Si  $Z \subset X_K$  est un sous-schéma fermé intègre 
avec image schématique  $\overline{Z} \subset X$, on voit que  
$Z_{K'} \subset X_{K'}$  est un sous-schéma fermé avec image schématique
  $\overline{Z}_{R'} \subset X_{R'}$. Le changement de base de  
$[Z]$ est  $[Z_{K'}]_r = [\overline{Z}_{K'}]_r$  par définition. Nous avons 
$$
sp_{X/R}([Z]) = [\overline{Z}_\kappa]_r
\quad\text{et}\quad
sp_{X'/R'}([\overline{Z}_{K'}]_r) = [(\overline{Z}_{R'})_{\kappa'}]_r
$$
 d'après le lemme \ref{lemma-specialization-module}. Comme  
$(\overline{Z}_{R'})_{\kappa'} = (\overline{Z}_\kappa)_{\kappa'}$
 , nous concluons. 
\end{proof} 

\begin{lemma} 
\label{lemma-specialization-flat-pullback} 
Soit  $R$  un anneau de valuation discrète avec corps des fractions $K$  et corps résiduel  
$\kappa$. Soit  $X$  un schéma localement de type fini sur  $R$. 
Soit  $f : X' \to X$  un morphisme qui est localement de type fini, plat 
et de dimension relative  $e$. Alors le diagramme 
$$
\xymatrix{
Z_{r + e}(X'_K) \ar[rr]_{sp_{X'/R}} & & Z_{r + e}(X'_\kappa) \\
Z_r(X_K) \ar[rr]^{sp_{X/R}} \ar[u] & & Z_r(X_\kappa) \ar[u]
}
$$
 est commutatif, où  $r \geq 0$  et les flèches verticales sont données par 
l'image inverse plate. 
\end{lemma} 

\begin{proof} 
Soit  $Z \subset X$  un sous-schéma fermé intègre dominant  $R$.
 Par la construction de  $sp_{X/R}$ , nous avons  $sp_{X/R}([Z_K]) = [Z_\kappa]_r$
  ce qui caractérise l'application de spécialisation. 
Posons  $Z' = f^{-1}(Z) = X' \times_X Z$.
 Comme  $R$  est un anneau de valuation,  $Z$  est plat sur  $R$. 
Par conséquent, $Z'$  est plat sur  $R$  et  
$sp_{X'/R}([Z'_K]_{r + e}) = [Z'_\kappa]_{r + e}$
  d'après le lemme \ref{lemma-specialization-closed}. 
D'après Homologie de Chow, lemme \ref{chow-lemma-pullback-coherent},
 on a $f_K^*[Z_K] = [Z'_K]_{r + e}$  et  
$f_\kappa^*[Z_\kappa]_r = [Z'_\kappa]_{r + e}$. Nous concluons. 
\end{proof} 

\begin{lemma} 
\label{lemma-specialization-proper-pushforward} 
Soit  $R$  un anneau de valuation discrète de corps des fractions $K$ et de corps résiduel  
$\kappa$. Soit  $f : X \to Y$  un morphisme propre de schémas localement 
de type fini sur  $R$. Alors le diagramme 
$$
\xymatrix{
Z_r(X_K) \ar[rr]_{sp_{X/R}} \ar[d] & & Z_r(X_\kappa) \ar[d] \\
Z_r(Y_K) \ar[rr]^{sp_{Y/R}} & & Z_r(Y_\kappa)
}
$$
 est commutatif, où  $r \geq 0$  et les flèches verticales sont données 
par l'image directe propre. 
\end{lemma} 

\begin{proof} 
Soit  $Z \subset X$  un sous-schéma fermé intègre dominant  $R$.
 Par la construction de  $sp_{X/R}$ , nous avons  $sp_{X/R}([Z_K]) = [Z_\kappa]_r$
  ce qui caractérise l'application de spécialisation. 
Posons  $Z' = f(Z) \subset Y$. Alors  $Z'$  est un sous-schéma fermé intègre 
de  $Y$  dominant  $R$. Ainsi $sp_{Y/R}([Z'_K]) = [Z'_\kappa]_r$. 

\medskip\noindent
 Nous pouvons considérer  $[Z]$  comme un élément de  $Z_{r + 1}(X)$. Par définition, 
nous avons  $f_*[Z] = 0$  si  $\dim(Z') < r + 1$  et  $f_*[Z] = d[Z']$
  si $Z \to Z'$  est génériquement fini de degré  $d$.
 Comme l'image directe propre commute avec l'image inverse plate par  $Y_K \to Y$
  (Homologie de Chow, lemme \ref{chow-lemma-flat-pullback-proper-pushforward}),
 on obtient respectivement  $f_{K, *}[Z_K] = 0$  ou  $f_{K, *}[Z_K] = d[Z'_K]$.
 Appliquons Homologie de Chow, lemme \ref{chow-lemma-closed-in-X-gysin}, 
au diagramme commutatif 
$$
\xymatrix{
X_\kappa \ar[d] \ar[r]_i & X \ar[d] \\
Y_\kappa \ar[r]^j & Y
}
$$
 Nous obtenons que  $f_{\kappa, *}[Z_\kappa]_r = 0$  ou  
$f_{\kappa, *}[Z_\kappa] = d[Z'_\kappa]_r$  car
  $i^*[Z] = [Z_k]_r$ et  $j^*[Z'] = [Z'_\kappa]_r$.
 En regroupant le tout, nous concluons. 
\end{proof} 









\section{Familles de cycles sur les fibres} 
\label{section-cycles-fibres} 

\noindent
 Soit $f : X \to S$ un morphisme de schémas localement de type fini et 
soit $r \geq 0$ un entier. Une 
{\it famille $\alpha$ de $r$-cycles sur les fibres de $X/S$} est une famille 
$$
\alpha = (\alpha_s)_{s \in S}
$$
 indexée par les points $s$ de $S$, où $\alpha_s \in Z_r(X_s)$
 est un $r$-cycle sur la fibre schématique $X_s$ de $f$ en $s$.
 Il existe diverses constructions que nous pouvons effectuer sur des familles de  
$r$-cycles sur les fibres. 

\medskip\noindent
 {\bf  Changement de base. }  Soit 
$$
\xymatrix{
X' \ar[r] \ar[d] & X \ar[d]^f \\
S' \ar[r]^g & S
}
$$
 un carré cartésien de morphismes de schémas avec  $f$  localement de type fini. 
Soit  $r \geq 0$  un entier. Étant donné une famille  $\alpha$  de $r$-cycles sur les
 fibres de  $X/S$ , nous définissons le {\it  changement de base }   $g^*\alpha$  de $\alpha$
  comme étant la famille 
$$
g^*\alpha = (\alpha'_{s'})_{s' \in S'}
$$
 où  $\alpha'_{s'} \in Z_r(X'_{s'})$  est le changement de 
base du cycle  $\alpha_s$  avec  $s = g(s')$ comme à 
la section  \ref{section-relative-fields}  via l'identification  
$X'_{s'} = X_s \times_{\Spec(\kappa(s))} \Spec(\kappa(s'))$
  des fibres schématiques. 

\medskip\noindent
 {\bf  Restriction.}  Soit  $f : X \to S$  un morphisme de schémas qui est localement
 de type fini. Soit  $r \geq 0$  un entier. Soient $U \subset X$  et  
$V \subset S$  des sous-schémas ouverts avec  $f(U) \subset V$. Étant donnée une famille  
$\alpha$  de  $r$-cycles sur les fibres de $X/S$ , nous pouvons définir la 
 {\it restriction}   $\alpha|_U$  de  $\alpha$  comme la 
famille de $r$-cycles sur les fibres de  $U/V$
$$
\alpha|_U = (\alpha_s|_{U_s})_{s \in V}
$$
 des restrictions aux fibres schématiques. 

\medskip\noindent
 {\bf Image inverse plate.}  Soit  $X \to S$  un morphisme de schémas qui est localement 
de type fini. Soient  $r, e \geq 0$  des entiers. Soit $f : X' \to X$  un 
morphisme plat, localement de type fini, et de dimension relative $e$. 
Étant donnée une famille  $\alpha$  de  $r$-cycles 
sur les fibres de  $X/S$ , nous définissons l'{\it image inverse plate}   $f^*\alpha$  de  $\alpha$
  comme la famille de $(r + e)$-cycles sur les fibres 
$$
f^*\alpha = (f_s^*\alpha_s)_{s \in S}
$$
 où  $f_s^*\alpha_s \in Z_{r + e}(X'_s)$  est l'image inverse 
plate du cycle  $\alpha_s$  dans  $Z_r(X_s)$  par le morphisme plat  
$f_s : X'_s \to X_s$  de dimension relative $e$
 induit sur les fibres schématiques. 

\medskip\noindent
 {\bf Image directe propre.}  Soit 
$$
\xymatrix{
X \ar[rr]_f \ar[rd] & & Y \ar[ld] \\
& S
}
$$
 un diagramme commutatif de morphismes de schémas avec  $X$  et $Y$
  localement de type fini sur  $S$  et  $f$  propre. Soit $r \geq 0$  un entier. Étant 
donné une famille  $\alpha$  de $r$-cycles sur les fibres de  $X/S$ , nous définissons 
l'{\it image directe propre}   $f_*\alpha$  de $\alpha$  comme étant la famille de  
$r$-cycles sur les fibres de $Y/S$  par 
$$
f_*\alpha = (f_{s, *}\alpha_s)_{s \in S}
$$
 où  $f_{s, *}\alpha_s \in Z_r(Y_s)$  est l'image directe propre
 du cycle  $\alpha_s$  dans  $Z_r(X_s)$ par le morphisme propre  
$f_s : X_s \to Y_s$  des fibres schématiques. 

\begin{lemma} 
\label{lemma-compatibilities} Les 
opérations précédentes vérifient les compatibilités suivantes : (1) le changement
 de base est fonctoriel ; (2) la 
restriction est la composée d'un changement de base et d'un cas particulier d'image inverse
 plate ; (3) 
l'image inverse plate commute au changement de base ; 
(4) l'image inverse plate est fonctorielle
 ; (5) l'image directe propre commute au changement de 
base ; (6) l'image directe propre est fonctorielle 
; (7) l'image directe propre commute à l'image inverse plate. 
\end{lemma} 

\begin{proof}
 Chacune de ces compatibilités découle directement des résultats 
correspondants prouvés dans le chapitre Homologie de Chow appliquée aux fibres
 sur $S$  des schémas en question. Nous omettons les énoncés précis et les 
preuves détaillées. Voici quelques références. Partie
 (1) : Homologie de Chow, lemme 
\ref{chow-lemma-compose-base-change}. Partie 
(2) : c'est immédiat. 
Partie (3) : Homologie de Chow, lemme 
\ref{chow-lemma-pullback-base-change-pullback}. Partie
 (4) : Homologie de Chow, lemme 
\ref{chow-lemma-compose-flat-pullback}. Partie
 (5) : Homologie de Chow, lemme 
\ref{chow-lemma-pullback-base-change-pushforward}. Partie
 (6) : Homologie de Chow, lemme 
\ref{chow-lemma-compose-pushforward}. Partie
 (7) : Homologie de Chow, lemme 
\ref{chow-lemma-flat-pullback-proper-pushforward}. 
\end{proof} 

\begin{example} 
\label{example-family-associated-module} 
Soit  $f : X \to S$  un morphisme de schémas qui est localement de type fini. 
Soit $r \geq 0$  un entier. Soit  $\mathcal{F}$  un  
$\mathcal{O}_X$-module quasi-cohérent de type fini. Pour  $s \in S$  notons  $\mathcal{F}_s$
  l'image inverse de $\mathcal{F}$ sur $X_s$. 
Supposons  $\dim(\text{Supp}(\mathcal{F}_s)) \leq r$  pour tout  $s \in S$.
 Alors nous pouvons associer à  $\mathcal{F}$  la famille  $[\mathcal{F}/X/S]_r$  de 
$r$-cycles sur les fibres de  $X/S$  définie par la formule 
$$
[\mathcal{F}/X/S]_r = ([\mathcal{F}_s]_r)_{s \in S}
$$
 où  $[\mathcal{F}_s]_r$  est donné par Homologie de Chow, définition  
\ref{chow-definition-cycle-associated-to-coherent-sheaf}. 
\end{example} 

\begin{lemma} 
\label{lemma-family-associated-module} La
 construction dans l'exemple  \ref{example-family-associated-module}  est 
compatible avec le changement de base, la restriction et l'image inverse plate. 
\end{lemma} 

\begin{proof} 
Voir Homologie de Chow, lemmes  
\ref{chow-lemma-pullback-coherent-base-change}  et  
\ref{chow-lemma-pullback-coherent}. 
\end{proof} 

\begin{example} 
\label{example-family-associated-closed} 
Soit  $f : X \to S$  un morphisme de schémas qui est localement de type fini. 
Soit $r \geq 0$  un entier. Soit  $Z \subset X$  un sous-schéma fermé. 
Pour  $s \in S$ , notons  $Z_s$  l'image réciproque de  $Z$  dans  $X_s$
  ou, de manière équivalente, la fibre schématique de  $Z$  en  $s$  vue 
comme un sous-schéma fermé de  $X_s$. 
Supposons  $\dim(Z_s) \leq r$  pour tout $s \in S$.
 Alors nous pouvons associer à  $Z$  la famille  $[Z/X/S]_r$
 de  $r$-cycles sur les fibres de  $X/S$  définie par la formule 
$$
[Z/X/S]_r = ([Z_s]_r)_{s \in S}
$$
 où  $[Z_s]_r$  est donné par
 Homologie de Chow, définition  
\ref{chow-definition-cycle-associated-to-closed-subscheme}. 
\end{example} 

\begin{lemma} 
\label{lemma-family-associated-closed} La
 construction dans l'exemple  \ref{example-family-associated-closed}  est 
compatible avec le changement de base, la restriction et l'image inverse plate. 
\end{lemma} 

\begin{proof} En
 prenant  $\mathcal{F} = (Z \to X)_*\mathcal{O}_Z$ , il s'agit d'un cas particulier du
 lemme  \ref{lemma-family-associated-module}. Voir 
Homologie de Chow, lemme \ref{chow-lemma-cycle-closed-coherent}. 
\end{proof} 

\begin{remark}[Support] 
\label{remark-supports-family} 
Soit $f : X \to S$ un morphisme de schémas qui est localement de type fini. 
Soit $r \geq 0$ un entier. Soit $\alpha$ une famille de $r$-cycles 
sur les fibres de $X/S$. Nous définissons le {\it support} de $\alpha$ comme 
$$
\text{Supp}(\alpha) =
\bigcup\nolimits_{s \in S} \text{Supp}(\alpha_s) \subset X
$$
 Ici,  $\text{Supp}(\alpha_s) \subset X_s$  est 
le support du cycle  $\alpha_s$, voir
 Homologie de Chow, définition  \ref{chow-definition-support-cycle}. 
Le support  $\text{Supp}(\alpha)$  est rarement un sous-ensemble fermé de  $X$. 
\end{remark} 

\begin{lemma} 
\label{lemma-support-family} La 
formation du support comme dans la remarque  \ref{remark-supports-family}  est 
compatible avec le changement de base, la restriction et l'image inverse plate. 
\end{lemma} 

\begin{proof} 
Démonstration omise. 
\end{proof} 

\begin{lemma} 
\label{lemma-coequalizer-dim-r} 
Soit  $f : X \to S$  un morphisme de schémas qui est localement de type fini. 
Soit $r \geq 0$  un entier. Soit  $g : S' \to S$  un morphisme de schémas 
surjectif. Posons  $S'' = S' \times_S S'$  et soient  $f' : X' \to S'$
  et  $f'' : X'' \to S''$  les changements de base de  $f$. 
Soit  $x \in X$  avec  $\text{trdeg}_{\kappa(f(x))}(\kappa(x)) = r$. 
\begin{enumerate} 
\item Il existe un point $x' \in X'$ se projetant sur $x$
  avec  $\text{trdeg}_{\kappa(f'(x'))}(\kappa(x')) = r$. 
\item Si  $x'_1, x'_2 \in X'$  sont tous deux comme dans (1), alors 
il existe un  $x'' \in X''$ avec  
$\text{trdeg}_{\kappa(f''(x''))}(\kappa(x'')) = r$  et  
$\text{pr}_i(x'') = x'_i$. 
\end{enumerate} 
\end{lemma} 

\begin{proof} La
 partie (1) résulte
 de Morphismes, lemme  \ref{morphisms-lemma-dimension-fibre-after-base-change}. 
Soit  $x'_1, x'_2$  comme dans (2). Alors, puisque  $X'' = X' \times_X X'$
 , nous voyons qu'il 
existe un point $x'' \in X''$ se projetant à la fois sur  $x'_1$  et  $x'_2$  (voir par
 exemple Descente, lemme \ref{descent-lemma-equiv-fibre-product}). On
 note  $s'' \in S''$,  $s'_i \in S'$ et  $s \in S$  les images 
de $x''$,  $x'_i$ et  $x$.
 Posons  $k = \kappa(s)$. Soit $Z \subset X_k$  le sous-schéma 
fermé intègre dont le point générique est $x$. Alors  $x'_i$
  est un point générique d'une composante irréductible de  $Z_{\kappa(s'_i)}$. 
Soit  $Z'' \subset Z_{\kappa(s'')}$  une composante irréductible
 contenant  $x''$. On note  $\xi'' \in Z''$  le point générique.
 Comme $\xi'' \leadsto x''$ , nous voyons que  $\xi''$  doit également se
 projeter sur $x'_i$ sous les deux projections. Par ailleurs, nous 
voyons que  $\text{trdeg}_{\kappa(s'')}(\kappa(\xi'')) = r$
  parce que c'est un point 
générique d'une composante irréductible du changement de base de $Z$. 
\end{proof} 

\begin{lemma} 
\label{lemma-descend-family} 
Soit  $f : X \to S$  un morphisme de schémas qui est localement de type fini. 
Soit $r \geq 0$  un entier. Soit  $g : S' \to S$  un morphisme de 
schémas et  $X' = S' \times_S X$. Supposons que pour tout  $s \in S$  il 
existe un point  $s' \in S'$  avec $g(s') = s$  et tel que  
$\kappa(s')/\kappa(s)$  est une extension séparable de corps. Alors 
\begin{enumerate} 
\item Si  $\alpha_1$  et  $\alpha_2$  sont des familles de  $r$-cycles sur les fibres de  $X/S$
  et si  $g^*\alpha_1 = g^*\alpha_2$, alors  $\alpha_1 = \alpha_2$. 
\item Étant donné une famille  $\alpha'$  de  $r$-cycles sur les fibres de $X'/S'$ , si  
$\text{pr}_1^*\alpha' = \text{pr}_2^*\alpha'$  en tant que familles de  
$r$-cycles sur les fibres de  $(S' \times_S S') \times_S X / (S' \times_S S')$,
 alors il existe une famille unique  $\alpha$  de $r$-cycles sur les fibres de  $X/S$
  telle que  $g^*\alpha = \alpha'$. 
\end{enumerate} 
\end{lemma} 

\begin{proof} 
La partie (1) découle de l'injectivité de l'application de changement de base étudiée à
 la section  \ref{section-relative-fields}. (Cet argument fonctionne 
tant que $S' \to S$  est surjectif.) 

\medskip\noindent
 Soit  $\alpha'$  comme dans (2). On note  
$\alpha'' = \text{pr}_1^*\alpha' = \text{pr}_2^*\alpha'$
  la valeur commune. 

\medskip\noindent
 Soit  $(X/S)^{(r)}$  l'ensemble des  $x \in X$  avec  
$\text{trdeg}_{\kappa(f(x))}(\kappa(x)) = r$
  et définissons de même  $(X'/S')^{(r)}$  et  $(X''/S'')^{(r)}$
 . En considérant les coefficients, nous pouvons considérer  $\alpha'$  et  $\alpha''$  comme des fonctions  
$\alpha' : (X'/S')^{(r)} \to \mathbf{Z}$ et  
$\alpha'' : (X''/S'')^{(r)} \to \mathbf{Z}$. 
Étant donnée une fonction 
$$
\varphi : (X/S)^{(r)} \to \mathbf{Z}
$$
 nous définissons  $g^*\varphi : (X'/S')^{(r)} \to \mathbf{Z}$  par analogie avec
 notre opération de changement de base. Plus précisément, supposons que  $x' \in (X'/S')^{(r)}$
  se projette sur $x \in X$,  $s' \in S'$ et $s \in Z$.
 On note  $Z' \subset X'_{s'}$  et  $Z \subset X_s$  les sous-schémas
 fermés intègres avec les points génériques  $x'$  et  $x$. Observons
 que $\dim(Z') = r$. Si  $\dim(Z) < r$, on pose  $(g^*\varphi)(x') = 0$. 
Si  $\dim(Z) = r$, alors  $Z'$  est une composante irréductible de  $Z_{s'}$  et 
possède donc une multiplicité  $m_{Z', Z_{s'}}$. Notons cette multiplicité $m(x', g)$.
 On définit alors 
$$
(g^*\varphi)(x') = m(x', g) \varphi(x)
$$
 Les coefficients  $m(x', g)$  sont toujours des entiers positifs
 (voir par exemple le lemme  \ref{lemma-multiplicities-field-extension}). De même, nous
 avons des applications de changement de base 
$$
\text{pr}_1^*, \text{pr}_2^* :
\text{Map}((X'/S')^{(r)}, \mathbf{Z})
\longrightarrow
\text{Map}((X''/S'')^{(r)}, \mathbf{Z})
$$
 Il découle de l'associativité du changement de base que nous avons  
$\text{pr}_1^* \circ g^* = \text{pr}_2^* \circ g^*$ (petit détail 
omis). Plus explicitement, en termes d'applications ensemblistes, cette 
égalité signifie simplement que pour  $x'' \in (X''/S'')^{(r)}$  nous avons 
$$
m(x'', \text{pr}_1) m(\text{pr}_1(x''), g) =
m(x'', \text{pr}_2) m(\text{pr}_2(x''), g)
$$
 pourvu que  $\text{pr}_1(x'')$  et  $\text{pr}_2(x'')$  soient 
dans  $(X''/S'')^{(r)}$.
 Le lemme \ref{lemma-coequalizer-dim-r}  et un argument 
élémentaire\footnote{Étant donné $x \in (X/S)^{(r)}$ , choisissez $x' \in (X'/S')^{(r)}$
 se projetant sur  $x$  et définissez $\alpha(x) = \alpha'(x')/m(x', g)$. Cela 
est bien défini par la formule et le lemme.}, 
appliqués à l'équation affichée précédente, 
montrent qu'il existe une application unique 
$$
\alpha : (X/S)^{(r)} \to \mathbf{Q}
$$
 telle que $g^*\alpha = \alpha'$. Pour terminer la preuve, il suffit 
de montrer que $\alpha$  prend des valeurs entières (petit détail omis:
 il faut voir que $\alpha$  détermine une somme localement finie
 sur chaque fibre, ce qui découle du fait correspondant pour  $\alpha'$). 
Étant donné un  $x \in (X/S)^{(r)}$  avec image $s \in S$
 , nous pouvons choisir un point  $s' \in S'$  tel que  $\kappa(s')/\kappa(s)$
  soit séparable. On peut ensuite choisir  $x' \in (X'/S')^{(r)}$  qui se projette
 sur $s$  et  $x$  et nous voyons que  $m(x', g) = 1$  parce que 
$Z_{s'}$  est réduit dans ce cas. Ainsi $\alpha(x) = \alpha'(x')$
 est un entier. 
\end{proof} 

\begin{lemma} 
\label{lemma-pullback-universally-bijective} 
Soit  $g : S' \to S$  un morphisme bijectif de schémas qui
 induit des isomorphismes des corps résiduels. 
Soit  $f : X \to S$  localement de type fini. Posons $X' = S' \times_S X$. 
Soit  $r \geq 0$. Alors le changement de base par  $g$ détermine une bijection
 entre le groupe des familles de  $r$-cycles sur les fibres de  $X/S$  et
 le groupe des familles de  $r$-cycles sur les fibres de  $X'/S'$. 
\end{lemma} 

\begin{proof} La 
démonstration est omise. 
\end{proof} 








\section{Cycles relatifs} 
\label{section-families-specialization} 

\noindent
 Voici la définition que nous utiliserons ; voir la section  \ref{section-compare}  
pour une comparaison avec les définitions de  \cite{SV}. 

\begin{definition} 
\label{definition-relative-cycles} 
Soit  $S$  un schéma localement noethérien. Soit  $f : X \to S$  un morphisme de 
schémas qui est localement de type fini. Soit  $r \geq 0$  un entier. 
Un {\it $r$-cycle relatif sur $X/S$} est une famille  $\alpha$  de  $r$-cycles 
sur les fibres de  $X/S$ telle que pour tout morphisme  $g : S' \to S$
  où  $S'$  est le spectre d'un anneau de valuation discrète, nous avons 
$$
sp_{X'/S'}(\alpha_\eta) = \alpha_0
$$
 où  $sp_{X'/S'}$  est défini à la section  \ref{section-specialization}  
et  $\alpha_\eta$  (resp. \ $\alpha_0$) est la valeur du changement de base  
$g^*\alpha$  de $\alpha$  au point générique (resp. fermé \ ) de  $S'$.
 Le groupe de tous les  $r$-cycles relatifs sur  $X/S$  est noté $z(X/S, r)$. 
\end{definition} 

\begin{lemma} 
\label{lemma-relative-cycle-functoriality} 
Soit  $\alpha$  un  $r$-cycle relatif sur  $X/S$  comme 
dans la définition  \ref{definition-relative-cycles}. Toute 
restriction, tout changement de base, toute image inverse plate ou toute image directe propre
 de $\alpha$ est alors un $r$-cycle relatif. 
\end{lemma} 

\begin{proof} Pour
 l'image inverse plate, on applique le lemme \ref{lemma-specialization-flat-pullback}. 
La restriction en est un cas particulier. Pour le changement de base, 
on utilise sa transitivité. Enfin, pour l'image 
directe propre, on applique le lemme \ref{lemma-specialization-proper-pushforward}. 
\end{proof} 

\begin{lemma} 
\label{lemma-relative-cycles-h-descent} 
Soit  $f : X \to S$  un morphisme de schémas. Supposons que  $S$  soit localement noethérien et
 que  $f$  soit localement de type fini. Soit  $r \geq 0$  un entier. Soit  $\alpha$
  une famille de  $r$-cycles sur les fibres de $X/S$. Soit  $\{g_i : S_i \to S\}$
  un recouvrement h (voir Compléments sur la platitude, définition  
\ref{flat-definition-h-covering}). Alors  $\alpha$  est un  $r$-cycle 
relatif si et seulement si chaque changement de base  $g_i^*\alpha$  est un  $r$-cycle relatif. 
\end{lemma} 

\begin{proof} 
Si  $\alpha$  est un  $r$-cycle relatif, alors chaque changement de base $g_i^*\alpha$  est
 un  $r$-cycle relatif d'après le lemme  \ref{lemma-relative-cycle-functoriality}. Supposons
 que chaque  $g_i^*\alpha$  soit un  $r$-cycle relatif.
 Soit $g : S' \to S$  un morphisme où  $S'$  est le spectre d'un anneau de 
valuation discrète. Après avoir remplacé  $S$  par  $S'$,  $X$  par $X' = X \times_S S'$, et  
$\alpha$  par  $\alpha' = g^*\alpha$, puis utilisé le fait que le changement de base d'un
 recouvrement h est un recouvrement h (Compléments sur la platitude, lemme \ref{flat-lemma-h}), 
nous revenons au problème étudié dans le paragraphe suivant. 

\medskip\noindent
 Supposons que  $S$  soit le spectre d'un anneau de valuation discrète avec le point 
fermé  $0$  et le point générique  $\eta$. Nous devons montrer que  
$sp_{X/S}(\alpha_\eta) = \alpha_0$. Puisqu'un recouvrement h est un recouvrement
 V (par définition), il existe un  $i$  et une spécialisation  $s' \leadsto s$
  des points de  $S_i$  avec  $g_i(s') = \eta$  et  $g_i(s) = 0$, voir
 Topologies, lemme \ref{topologies-lemma-refine-qcqs-V}. D'après
 Propriétés, lemme  \ref{properties-lemma-locally-Noetherian-specialization-dvr} , nous 
pouvons trouver un morphisme $h : S' \to S_i$ du 
spectre  $S'$  d'un anneau de valuation discrète qui envoie
 le point générique  $\eta'$  sur  $s'$  et envoie
 le point fermé  $0'$  sur  $s$. Notons  $\alpha' = h^*g_i^*\alpha$.
 Par hypothèse, nous avons  $sp_{X'/S'}(\alpha'_{\eta'}) = \alpha'_{0'}$. 
Puisque  $g = g_i \circ h : S' \to S$  est le morphisme de schémas 
induit par une extension d'anneaux de valuation discrète, nous en déduisons que 
$sp_{X/S}$  et  $sp_{X'/S'}$  sont compatibles avec les applications de changement 
de base sur les fibres, voir le lemme  \ref{lemma-specialization-extension}. 
Nous concluons que $sp_{X/S}(\alpha_\eta) = \alpha_0$  car 
l'application de changement de base  $Z_r(X_0) \to Z_r(X'_{0'})$  est injective 
comme on l'a vu à la section  \ref{section-relative-fields}. 
\end{proof} 

\begin{lemma} 
\label{lemma-families-specialization-fppf-descent} 
Soit  $f : X \to S$  un morphisme de schémas. Supposons que  $S$  soit localement noethérien et
 que  $f$  soit localement de type fini. Soient  $r, e \geq 0$  des entiers. 
Soit  $\alpha$  une famille de  $r$-cycles sur les fibres de $X/S$. 
Soit  $\{f_i : X_i \to X\}$  une famille de morphismes
 plats, localement de type fini, et de dimension relative  $e$, qui 
est conjointement surjective. Alors  $\alpha$  est un  $r$-cycle relatif si et seulement si 
chaque cycle $f_i^*\alpha$ obtenu par image inverse plate  est un  $r$-cycle relatif. 
\end{lemma} 

\begin{proof} 
Si  $\alpha$  est un  $r$-cycle relatif, alors chaque  $f_i^*\alpha$ est à 
son tour un $r$-cycle relatif d'après le lemme  \ref{lemma-relative-cycle-functoriality}. Supposons
 que chaque  $f_i^*\alpha$  soit un  $r$-cycle relatif.
 Soit  $g : S' \to S$  un morphisme où  $S'$  est le spectre d'un anneau de 
valuation discrète. Après avoir remplacé  $S$  par  $S'$,  $X$  par  $X' = X \times_S S'$, et 
$\alpha$  par  $\alpha' = g^*\alpha$
 , nous réduisons au problème étudié dans le paragraphe suivant. 

\medskip\noindent
 Supposons que  $S$  soit le spectre d'un anneau de valuation discrète avec point 
fermé  $0$  et point générique  $\eta$. Nous devons montrer que 
$sp_{X/S}(\alpha_\eta) = \alpha_0$. Notons  $f_{i, 0} : X_{i, 0} \to X_0$
  le changement de base de  $f_i$  au point fermé de  $S$. De même 
pour  $f_{i, \eta}$.
 Observons que 
$$
f_{i, 0}^*sp_{X/S}(\alpha_\eta) =
sp_{X_i/S}(f_{i, \eta}^*\alpha_\eta) = f_{i, 0}^*\alpha_0
$$
 En effet, la première égalité résulte du 
lemme  \ref{lemma-specialization-flat-pullback}  et
 la seconde par hypothèse. Puisque la famille de morphismes  
$f_{i, 0}^* : Z_r(X_0) \to Z_r(X_{i, 0})$  est 
conjointement injective (puisque  $f_{i, 0}$  est conjointement
 surjective), nous concluons ce que nous voulons. 
\end{proof} 

\begin{lemma} 
\label{lemma-check-after-closed} 
Soit  $S$  un schéma localement noethérien. Soit  $i : X \to Y$  une immersion fermée 
de schémas localement de type fini sur  $S$. Soit  $r \geq 0$. 
Soit  $\alpha$  une famille de  $r$-cycles sur les fibres de $X/S$. 
Alors  $\alpha$  est un  $r$-cycle relatif sur $X/S$  si et seulement si  
$i_*\alpha$  est un  $r$-cycle relatif sur $Y/S$. 
\end{lemma} 

\begin{proof} 
Puisque le changement de base commute avec  $i_*$  (lemme \ref{lemma-compatibilities}),
 il suffit de montrer ceci : si  $S$  est le spectre d'un anneau de 
valuation discrète avec point générique  $\eta$  et point fermé  $0$, 
alors  $sp_{X/S}(\alpha_\eta) = \alpha_0$  si et seulement si  
$sp_{Y/S}(i_{\eta, *}\alpha_\eta) = i_{0, *}\alpha_0$. 
Cela résulte de l'injectivité de  $i_{0, *} : Z_r(X_0) \to Z_r(Y_0)$
 et de l'égalité  $i_{0, *}sp_{X/S}(\alpha_\eta) =
sp_{Y/S}(i_{\eta, *}\alpha_\eta)$  d'après
 le lemme  \ref{lemma-specialization-proper-pushforward}. 
\end{proof} 

\noindent
 Le lemme suivant sera renforcé dans le 
lemme  \ref{lemma-relative-cycles-equal}. 

\begin{lemma} 
\label{lemma-uniqueness-extension} 
Soit  $f : X \to S$  un morphisme de schémas. Supposons que  $S$  soit localement noethérien
 et  $f$  localement de type fini. Soit  $r \geq 0$. Soient  $\alpha$  et  $\beta$
  des  $r$-cycles relatifs sur $X/S$. Les énoncés suivants sont équivalents 
\begin{enumerate} 
\item $\alpha = \beta$, et 
\item $\alpha_\eta = \beta_\eta$  pour tout point générique  $\eta \in S$
  d'une composante irréductible de  $S$. 
\end{enumerate} 
\end{lemma} 

\begin{proof}
 L'implication (1)  $\Rightarrow$  (2) est immédiate. 
Supposons (2). Pour chaque $s \in S$  nous pouvons trouver un point $\eta$ comme en (2)
 qui se spécialise en  $s$. 
D'après Propriétés, lemme  \ref{properties-lemma-locally-Noetherian-specialization-dvr} , nous 
pouvons trouver un morphisme  $g : S' \to S$  depuis le
 spectre  $S'$  d'un anneau de valuation discrète qui envoie
 le point générique  $\eta'$  sur  $\eta$  et qui 
envoie le point fermé  $0$  sur  $s$. Alors  $\alpha_s$  et  $\beta_s$
 sont des éléments de  $Z_r(X_s)$  dont les images par changement de base 
sont égales dans  $Z_r(X_{0'})$, à savoir  $sp_{X_{S'}/S'}(\alpha_{\eta'})$
  où  $\alpha_{\eta'}$  est le changement de base de  $\alpha_\eta$. 
Comme l'application de changement de base $Z_r(X_s) \to Z_r(X_{0'})$  est injective 
comme on l'a vu à la section  \ref{section-relative-fields} , nous
 en concluons  $\alpha_s = \beta_s$. 
\end{proof} 

\begin{lemma} 
\label{lemma-family-associated-module-specialization} 
Dans la situation de l'exemple  \ref{example-family-associated-module} , supposons
 que  $S$  soit localement noethérien et que  
$\mathcal{F}$  soit plat sur  $S$  en dimensions $\geq r$
  (Compléments sur la platitude, définition  \ref{flat-definition-flat-dimension-n}). 
Alors  $[\mathcal{F}/X/S]_r$ est un  $r$-cycle relatif sur  $X/S$. 
\end{lemma} 

\begin{proof} 
D'après Compléments sur la platitude, lemme \ref{flat-lemma-pre-flat-dimension-n} 
, l'hypothèse sur  $\mathcal{F}$ est conservée par tout changement de base. 
De plus, la formation de  $[\mathcal{F}/X/S]_r$ est compatible avec tout changement
 de base d'après le lemme  \ref{lemma-family-associated-module}. 
Puisque la condition d'être compatible avec les spécialisations est
 vérifiée après un changement de base par le spectre d'un anneau de valuation 
discrète, cela nous ramène au cas où  $S$  est le spectre d'un anneau de valuation. 
Dans ce cas, l'ensemble  
$U = \{x \in X \mid \mathcal{F}\text{ flat at }x\text{ over }S\}$
  est ouvert dans  $X$  
d'après Compléments sur la platitude, lemme \ref{flat-lemma-finite-type-flat}. 
Puisque le complément de  $U$  dans $X$  a des fibres de dimension  $< r$  sur  
$S$  par hypothèse, nous voyons que la restriction à l'ouvert  
$U \subset X$ induit un isomorphisme sur les groupes de  $r$-cycles 
sur les fibres après tout changement de base, compatible avec les applications de
 spécialisation et avec la formation du cycle relatif associé à  $\mathcal{F}$.
 Ainsi, il suffit de montrer la compatibilité avec les
 spécialisations pour  $[\mathcal{F}|_U / U /S]_r$.
 Puisque $\mathcal{F}|_U$  est plat sur  $S$, cela découle du 
lemme  \ref{lemma-specialization-module}  et des définitions. 
\end{proof} 

\begin{lemma} 
\label{lemma-family-associated-closed-specialization} 
Dans la situation de l'exemple  \ref{example-family-associated-closed} , supposons
 que $S$ soit localement noethérien et que  $Z$  soit plat sur  $S$  en dimensions 
$\geq r$. Alors  $[Z/X/S]_r$  est un  $r$-cycle relatif sur $X/S$. 
\end{lemma} 

\begin{proof}
 L'hypothèse signifie que  $\mathcal{O}_Z$  est plat sur  $S$  
en dimensions $\geq r$. Ainsi, en appliquant 
le lemme  \ref{lemma-family-associated-module-specialization}  
avec  $\mathcal{F} = (Z \to X)_*\mathcal{O}_Z$ , nous en concluons. 
\end{proof} 

\noindent
 Soit  $S$  un schéma localement noethérien. Soit  $f : X \to S$  un morphisme
 de type fini. Soit  $r \geq 0$. Notons  $Hilb(X/S, r)$
  l'ensemble des sous-schémas fermés  $Z \subset X$  tels que  $Z \to S$  soit 
plat et de dimension relative $\leq r$. D'après
 le lemme  \ref{lemma-family-associated-closed-specialization} , pour chaque  
$Z \in Hilb(X/S, r)$ , nous avons un élément  $[Z/X/S]_r \in z(X/S, r)$.
 Ainsi, nous obtenons un homomorphisme de groupe 
\begin{equation} 
\label{equation-cycle-classes} 
\text{groupe abélien libre sur }Hilb(X/S, r) \longrightarrow z(X/S, r) 
\end{equation} 
envoyant  $\sum n_i[Z_i]$  vers  $\sum n_i[Z_i/X/S]_r$.
 Une propriété essentielle des $r$-cycles relatifs est qu'ils sont localement 
(sur  $X$  et $S$  dans des topologies appropriées) dans l'image de cette application. 

\begin{lemma} 
\label{lemma-get-cycles} 
Soit  $f : X \to S$  un morphisme de type fini entre schémas, où $S$ est noethérien. 
Soit  $r \geq 0$. Soit  $\alpha$  un  $r$-cycle relatif sur $X/S$. Alors il existe un 
morphisme propre, complètement décomposé 
(Compléments sur les morphismes, définition  \ref{more-morphisms-definition-cd-morphism})
  $g : S' \to S$  tel que  $g^*\alpha$ soit dans l'image de 
(\ref{equation-cycle-classes}). 
\end{lemma} 

\begin{proof} 
Par induction noethérienne, on peut supposer que le résultat est vrai pour l'image inverse de 
$\alpha$  par toute immersion fermée  $g : S' \to S$  qui n'est pas un isomorphisme. 

\medskip\noindent
 Soit  $S_1 \subset S$  une composante irréductible (vue comme un sous-schéma fermé
 intègre). Soit  $S_2 \subset S$  l'adhérence du complément de  $S'$
  (vue comme un sous-schéma fermé réduit). Si  $S_2 \not = \emptyset$, alors
 le résultat est vrai pour l'image inverse de $\alpha$  par  $S_1 \to S$  et  $S_2 \to S$. 
Si  $g_1 : S'_1 \to S_1$ et  $g_2 : S'_2 \to S_2$
  sont les morphismes propres complètement décomposés correspondants, 
alors  $S' = S'_1 \amalg S'_2 \to S$
  est un morphisme propre complètement décomposé et nous 
voyons que le résultat est valable pour $S$\footnote{Plus précisément, tout 
sous-schéma fermé de $S'_1 \times_S X$ plat et de dimension relative  
$\leq r$  sur  $S'_1$  peut être considéré comme un sous-schéma fermé de  $S' \times_S X$
  plat et de dimension relative  $\leq r$  sur  $S'$.} 
. Ainsi, nous pouvons supposer que  $S' \to S$  est bijectif
 et nous nous ramenons au cas décrit dans le paragraphe suivant. 

\medskip\noindent
 Supposons que  $S$  soit intègre. Soit  $\eta \in S$  le point générique 
et soit  $K = \kappa(\eta)$  le corps des fonctions de  $S$. 
Alors  $\alpha_\eta$  est un  $r$-cycle sur  $X_K$. 
Écrivons  $\alpha_\eta = \sum n_i[Y_i]$.
 En prenant l'adhérence de  $Y_i$ , nous obtenons des sous-schémas fermés intègres 
$Z_i \subset X$  dont le changement de base vers  $\eta$  est  $Y_i$.
 Par platitude générique (par exemple Morphismes,
 proposition  \ref{morphisms-proposition-generic-flatness}), nous
 voyons que  $Z_i$  est plat sur un ouvert non vide  $U$  de  $S$  pour chaque  $i$. 
En appliquant Compléments sur la platitude, lemme \ref{flat-lemma-flat-after-blowing-up} , nous pouvons 
trouver un éclatement  $U$-admissible  $g : S' \to S$
 tel que la transformée stricte  $Z'_i \subset X_{S'}$
  de  $Z_i$  soit plate sur $S'$. Alors  $\beta = \sum n_i[Z'_i/X_{S'}/S']_r$
  est dans l'image de (\ref{equation-cycle-classes}) et $\beta = g^*\alpha$
  par le lemme \ref{lemma-uniqueness-extension}. 

\medskip\noindent
 Cependant, cela ne termine pas la preuve car  $S' \to S$  peut ne pas 
être complètement décomposé. Cela se corrige facilement : en notant  $T \subset S$
  le complément de  $U$  (vu comme un sous-schéma fermé), par induction 
noethérienne nous pouvons trouver un morphisme propre complètement décomposé 
$T' \to T$  tel que  $(T' \to S)^*\alpha$
  soit dans l'image de (\ref{equation-cycle-classes}). Alors 
$S' \amalg T' \to S$  convient. 
\end{proof} 

\begin{lemma} 
\label{lemma-get-cycles-dvr} 
Soit  $f : X \to S$  un morphisme de schémas de type fini avec  $S$
  le spectre d'un anneau de valuation discrète. Soit  $r \geq 0$.
 Alors (\ref{equation-cycle-classes}) est surjectif. 
\end{lemma} 

\begin{proof}
 Cela découle bien sûr du lemme \ref{lemma-get-cycles}, mais nous pouvons 
également le voir directement comme suit. Supposons que  $\alpha$  soit un  $r$-cycle relatif
 sur  $X/S$. Écrivons  $\alpha_\eta = \sum n_i[Z_i]$  (la somme est finie). Notons 
$\overline{Z}_i \subset X$  l'adhérence de  $Z_i$  comme 
dans la section \ref{section-specialization}. Alors 
$\alpha = \sum n_i[\overline{Z}_i/X/S]$. 
\end{proof} 

\begin{lemma} 
\label{lemma-relative-cycle-smooth} 
Soit  $f : X \to S$  un morphisme de schémas. Soit  $r \geq 0$. Supposons que $S$
 soit localement noethérien et que $f$ soit lisse de dimension relative $r$. Soit  
$\alpha \in z(X/S, r)$. Alors le support de  $\alpha$  est ouvert et fermé dans  $X$
  (voir la preuve pour un résultat plus précis). 
\end{lemma} 

\begin{proof} 
Soit  $x \in X$  d'image  $s \in S$. Puisque  $f$  est lisse, il existe une
 unique composante irréductible  $Z(x)$  de  $X_s$  qui contient  $x$. 
Alors  $\dim(Z(x)) = r$. 
Soit  $n_x$  le coefficient de $Z(x)$  dans le cycle  $\alpha_s$.
 Nous montrerons que la fonction $x \mapsto n_x$  est localement constante sur  $X$. 

\medskip\noindent
 Soit  $g : S' \to S$  un morphisme de schémas localement noethériens. 
Soit $X'$  le changement de base de  $X$  et soit  $\alpha' = g^*\alpha$
  le changement de base de  $\alpha$. Soit  $x' \in X'$  se projetant sur $s' \in S'$,  
$x \in X$ et  $s \in S$. Nous affirmons que $n_{x'} = n_x$. En effet, puisque  $Z(x)$
  est lisse sur  $\kappa(s)$ , nous voyons que  
$Z(x) \times_{\Spec(\kappa(s))} \Spec(\kappa(s'))$
  est réduit. Puisque  $Z(x')$  est une composante irréductible
 de ce schéma, nous voyons que le coefficient  $n_{x'}$  de  $Z(x')$
  dans  $\alpha'_{s'}$  est le même que le coefficient  $n_x$  de  $Z(x)$
  dans  $\alpha_s$  par la définition du changement de base dans 
la section \ref{section-relative-fields}, 
ce qui prouve l'affirmation. 

\medskip\noindent
 Puisque  $X$  est localement noethérien, pour montrer que  $x \mapsto n_x$
  est localement constant, il suffit de montrer: si  $x' \leadsto x$
  est une spécialisation dans  $X$, alors  $n_{x'} = n_x$.
 Choisissons un morphisme  $S' \to X$  où $S'$  est le spectre d'un anneau de 
valuation discrète envoyant le point générique  $\eta$  vers  $x'$  et le point fermé  
$0$  vers $x$. Voir Propriétés, lemme 
\ref{properties-lemma-locally-Noetherian-specialization-dvr}. Ensuite,
 le changement de base $X' \to S'$  de  $f$  par  $S' \to S$
  a une section  $\sigma : S' \to X'$  telle que $\sigma(\eta) \leadsto \sigma(0)$
  est une spécialisation de points de  $X'$  envoyés vers $x' \leadsto x$  dans  $X$.
 Ainsi, nous nous réduisons à l'affirmation du paragraphe suivant. 

\medskip\noindent
 Soit  $S$  le spectre d'un anneau de valuation discrète, de 
point générique  $\eta$  et de point fermé  $0$, et supposons donnée une section 
$\sigma : S  \to X$. Nous affirmons que $n_{\sigma(\eta)} = n_{\sigma(0)}$.
 D'après la discussion dans Compléments sur les morphismes, section 
\ref{more-morphisms-section-connected-components}, et en particulier Compléments
 sur les morphismes, lemme \ref{more-morphisms-lemma-connected-along-section-open} , après
 avoir remplacé  $X$  par un sous-schéma ouvert, nous pouvons supposer que 
les fibres de  $X \to S$  sont connexes. Comme ces fibres sont
 lisses, elles sont irréductibles. Alors nous voyons que  $\alpha_\eta = n[X_\eta]$
  avec $n = n_{\sigma(\eta)}$  et la relation  $sp_{X/S}(\alpha_\eta) = \alpha_0$
  implique  $\alpha_0 = n[X_0]$, c'est-à-dire $n_{\sigma(0)} = n$, comme voulu. 
\end{proof} 

\begin{lemma} 
\label{lemma-relative-cycles-equal} 
Soit  $f : X \to S$  un morphisme de schémas. Supposons que  $S$  soit localement noethérien
 et  $f$  localement de type fini. Soit  $r \geq 0$ et soient  
$\alpha, \beta \in z(X/S, r)$. L'ensemble  $E = \{s \in S : \alpha_s = \beta_s\}$
  est fermé dans  $S$. 
\end{lemma} 

\begin{proof} 
La question est locale sur  $S$, donc nous pouvons supposer que $S$  est affine. 
Soit  $X = \bigcup U_i$  un recouvrement ouvert affine. Soit 
$E_i = \{s \in S : \alpha_s|_{U_{i, s}} = \beta_s|_{U_{i, s}}\}$. 
Alors  $E = \bigcap E_i$. Par conséquent, il suffit de prouver le lemme pour  
$U_i \to S$  et la restriction de  $\alpha$  et  $\beta$  à $U_i$.
 Cela nous ramène au cas discuté dans le paragraphe suivant. 

\medskip\noindent
 Supposons que  $X$  et  $S$  soient quasi-compacts. 
Posons $\gamma = \alpha - \beta$. Alors  $E = \{s \in S : \gamma_s = 0\}$. 
D'après le lemme \ref{lemma-family-associated-closed-specialization}, il 
existe une famille finie de morphismes propres conjointement surjective  
$\{g_i : S_i \to S\}$  telle que  $g_i^*\gamma$  soit dans l'image
 de (\ref{equation-cycle-classes}). On observe que  $E_i = g_i^{-1}(E)$ est
 l'ensemble des points  $t \in S_i$  tels que  $(g_i^*\gamma)_t = 0$. 
Si  $E_i$  est fermé pour tous les  $i$, alors  $E = \bigcup g_i(E_i)$
  est également fermé. Cela nous ramène au cas discuté dans le paragraphe
 suivant. 

\medskip\noindent
 Supposons que  $X$  et  $S$  soient quasi-compacts et  $\gamma = \sum n_i[Z_i/X/S]_r$
 pour un nombre fini de sous-schémas fermés  $Z_i \subset X$
  plats et de dimension relative  $\leq r$  sur  $S$. 
Posons  $X' = \bigcup Z_i$  (union schématique). 
Alors  $i : X' \to X$  est une immersion fermée et  $X'$
  a une dimension relative  $\leq r$  sur  $S$. De plus  
$\gamma = i_*\gamma'$  où  $\gamma' = \sum n_i[Z_i/X'/S]_r$.
 Puisque  $E = E' = \{s \in S : \gamma'_s = 0\}$
 , nous nous réduisons au cas discuté dans le paragraphe suivant. 

\medskip\noindent
 Supposons que  $X$  ait une dimension relative  $\leq r$  sur $S$. 
Soit  $s \in S$, avec  $s \not \in E$. Nous allons montrer qu'il existe 
un voisinage ouvert  $V \subset S$  de  $s$  tel que  $E \cap V$  soit vide.
 L'hypothèse  $s \not \in E$  signifie qu'il existe un sous-schéma fermé
 intègre $Z \subset X_s$  de dimension  $r$  tel que le coefficient  
$n$  de $[Z]$  dans  $\gamma_s$  soit non nul. Soit  $x \in Z$  le point
 générique. Puisque  $\dim(Z) = r$ , nous voyons que  $x$  est un point 
générique d'une composante irréductible (à savoir  $Z$) de  $X_s$.
 Ainsi, après avoir remplacé  $X$  par un voisinage ouvert de  $x$,
 nous pouvons supposer que  $Z$  est la seule composante irréductible de  $X_s$.
 En particulier, nous avons  $\gamma_s = n[Z]$. 

\medskip\noindent
 À ce stade, nous appliquons Compléments sur les morphismes, lemme 
\ref{more-morphisms-lemma-local-structure-finite-type}, et 
nous obtenons un diagramme 
$$
\xymatrix{
X \ar[dd] & X' \ar[l]^g \ar[d]^\pi & x \ar@{|->}[dd] &
x' \ar@{|->}[l]  \ar@{|->}[d] \\
& Y \ar[d]^h & & y \ar@{|->}[d] \\
S \ar@{=}[r] & S & s & s \ar@{=}[l]
}
$$
 avec toutes les propriétés qui y sont énumérées. Soit  $\gamma' = g^*\gamma$
  l'image inverse plate. Notons que  $E \subset E' = \{s \in S: \gamma'_s = 0\}$
  et que  $s \not \in E'$  parce que le coefficient de $Z'$  dans  $\gamma'_s$
  est non nul, où  $Z' \subset X'_s$  est l'adhérence de $x'$.
 De même, posons  $\gamma'' = \pi_*\gamma'$. Alors nous avons  
$E' \subset E'' = \{s \in S: \gamma''_s = 0\}$  et $s \not \in E''$
  parce que le coefficient de  $Z''$  dans  $\gamma''_s$  est non nul, où  
$Z'' \subset Y_s$  est l'adhérence de  $y$. D'après
 le lemme \ref{lemma-relative-cycle-smooth}  et l'ouverture de  $Y \to S$
 , nous voyons qu'un voisinage ouvert de $s$  est disjoint de  $E''$
  et la démonstration est complète. 
\end{proof} 

\begin{lemma} 
\label{lemma-descend-through-limit} 
Soit  $S = \lim_{i \in I} S_i$  la limite d'un système inverse dirigé de schémas
 noethériens avec des morphismes de transition affines. 
Soit  $0 \in I$  et soit  $X_0 \to S_0$ un morphisme de schémas de type fini. 
Pour  $i \geq 0$ , posons  $X_i = S_i \times_{S_0} X_0$  et posons  
$X = S \times_{S_0} X_0$. Si  $S$  est également noethérien, alors 
$$
z(X/S, r) = \colim_{i \geq 0} z(X_i/S_i, r)
$$
 où les applications de transition sont données par changement de base des 
$r$-cycles relatifs. 
\end{lemma} 

\begin{proof} 
Supposons que  $i \geq 0$  et  $\alpha_i, \beta_i \in z(X_i/S_i, r)$
  aient la même image dans $z(X/S, r)$. Alors l'image de  $S \to S_i$
  est contenue dans le sous-ensemble fermé $E \subset S_i$  du
 lemme \ref{lemma-relative-cycles-equal}. Par conséquent,
 pour un certain  $j \geq i$ , l'image du morphisme  $S_j \to S_i$
  est contenue dans  $E$ ; voir Limites, lemme 
\ref{limits-lemma-limit-contained-in-constructible}. Il 
s'ensuit que les changements de base de  
$\alpha_i$  et $\beta_i$  vers  $S_j$  coïncident. Ainsi, l'application est injective. 

\medskip\noindent
 Soit  $\alpha \in z(X/S, r)$. En appliquant le lemme \ref{lemma-get-cycles}, on 
obtient un morphisme propre complètement décomposé  $g : S' \to S$  tel que  $g^*\alpha$
  soit dans l'image de (\ref{equation-cycle-classes}). 
Posons  $X' = S' \times_S X$. Nous écrivons  $g^*\alpha = \sum n_a [Z_a/X'/S']_r$
  pour certains sous-schémas fermés  $Z_a \subset X'$  plats et 
de dimension relative  $\leq r$  sur $S'$. 

\medskip\noindent
 Nous appliquons maintenant les résultats
 de Limites, section \ref{limits-section-descending-relative} et suivantes. Nous 
pouvons trouver un  $i \geq 0$  tel qu'il existe 
\begin{enumerate} 
\item un morphisme propre complètement décomposé  $g_i : S'_i \to S_i$
  dont le changement de base à  $S$  est  $g : S' \to S$, 
\item en posant  $X'_i = S'_i \times_{S_i} X_i$
 , des sous-schémas fermés  $Z_{ai} \subset X'_i$  plats et de
 dimension relative  $\leq r$  sur  $S'_i$  dont le changement de base en $S'$
  soit  $Z_a$. 
\end{enumerate} 
Pour ce faire, on utilise Limites, lemmes 
\ref{limits-lemma-descend-finite-presentation}, 
\ref{limits-lemma-descend-closed-immersion-finite-presentation}, 
\ref{limits-lemma-descend-flat-finite-presentation}, 
\ref{limits-lemma-eventually-proper} et 
\ref{limits-lemma-limit-dimension}, ainsi que
 Compléments sur les morphismes, lemme 
\ref{more-morphisms-lemma-descend-cd}. 
Considérons  
$\alpha'_i = \sum n_a [Z_{ai}/X'_i/S'_i]_r \in z(X'_i/S'_i, r)$.
 L'image de  $\alpha'_i$  dans  $z(X'/S', r)$ coïncide par construction avec le changement de base  
$g^*\alpha$. 

\medskip\noindent
 Posons  $S''_i = S'_i \times_{S_i} S'_i$  et  $X''_i = S''_i \times_{S_i} X_i$
  et posons  $S'' = S' \times_S S'$  et  $X'' = S'' \times_S X$.
 Nous notons  $\text{pr}_1, \text{pr}_2 : S'' \to S'$  et  
$\text{pr}_1, \text{pr}_2 : S''_i \to S'_i$  les projections. Les
 deux changements de base $\text{pr}_1^*\alpha'_i$  et  $\text{pr}_1^*\alpha'_i$
  ont la même image dans  $z(X''/S'', r)$  parce que  
$\text{pr}_1^*g^*\alpha = \text{pr}_1^*g^*\alpha$.
 Ainsi, quitte à remplacer  $i$  par un indice plus grand, nous pouvons supposer que  
$\text{pr}_1^*\alpha'_i = \text{pr}_1^*\alpha'_i$
  d'après le premier paragraphe de la démonstration. 
Par le lemme \ref{lemma-descend-family} , 
nous obtenons une famille unique  $\alpha_i$
  de  $r$-cycles sur les fibres de  $X_i/S_i$
  avec  $g_i^*\alpha_i = \alpha'_i$  (cela utilise le fait que $S'_i \to S_i$
  est complètement décomposé). Par
 le lemme \ref{lemma-relative-cycles-h-descent}, nous
 voyons que  $\alpha_i \in z(X_i/S_i, r)$.
 L'unicité dans le lemme \ref{lemma-descend-family} implique que 
l'image de $\alpha_i$  dans  $z(X/S, r)$  est  $\alpha$  et la démonstration est complète. 
\end{proof} 

\begin{lemma} 
\label{lemma-thickening} 
Soit  $S$  un schéma localement noethérien. Soit  $i : X \to X'$  un épaississement
 de schémas localement de type fini sur  $S$. Soit $r \geq 0$. 
Alors  $i_* : z(X/S, r) \to z(X'/S, r)$  est une bijection. 
\end{lemma} 

\begin{proof} 
Puisque  $i_s : X_s \to X'_s$  est un épaississement, il est clair que  $i_*$  induit 
une bijection entre les
 familles de  $r$-cycles sur les fibres de $X/S$  et 
les familles de  $r$-cycles sur les fibres de $X'/S$. 
De plus, étant donnée une famille  $\alpha$  de $r$-cycles sur les fibres de  $X/S$
 . De plus, $\alpha \in z(X/S, r) \Leftrightarrow i_*\alpha \in z(X'/S, r)$
  par le lemme \ref{lemma-check-after-closed}. Le lemme s'ensuit. 
\end{proof} 

\begin{lemma} 
\label{lemma-extend-to-larger} 
Soit  $S$  un schéma localement noethérien. Soit  $X$  un schéma localement
 de type fini sur  $S$. Soit  $r \geq 0$. Soit  $U \subset X$  un ouvert
 tel que  $X \setminus U$  ait une dimension relative  $< r$  sur $S$, c'est-à-dire  
$\dim(X_s \setminus U_s) < r$  pour tout  $s \in S$. Alors
 la restriction définit une bijection  $z(X/S, r) \to z(U/S, r)$. 
\end{lemma} 

\begin{proof} Étant
 donné que  $Z_r(X_s) \to Z_r(U_s)$  est une bijection par l'hypothèse de dimension, 
nous voyons que la restriction induit une bijection entre les
 familles de $r$-cycles sur les fibres de  $X/S$  et 
les familles de $r$-cycles sur les fibres de  $U/S$.
 Ces applications de restriction  $Z_r(X_s) \to Z_r(U_s)$  sont
 compatibles avec le changement de base et avec les spécialisations ; voir les
 lemmes \ref{lemma-compatibilities} et \ref{lemma-specialization-flat-pullback}. Le 
lemme s'en déduit facilement ; la démonstration est omise. 
\end{proof} 

\begin{lemma} 
\label{lemma-seminormalize-base} 
Soit  $g : S' \to S$  un homéomorphisme universel de schémas localement noethériens 
qui induit des isomorphismes des corps résiduels. Soit  $f : X \to S$  localement de 
type fini. Posons  $X' = S' \times_S X$. Soit  $r \geq 0$. Alors le changement de base par $g$
  détermine une bijection  $z(X/S, r) \to z(X'/S', r)$. 
\end{lemma} 

\begin{proof} 
D'après le lemme \ref{lemma-pullback-universally-bijective}, nous avons une bijection
 entre le groupe des familles de  $r$-cycles sur les fibres de  $X/S$  et
 le groupe des familles de  $r$-cycles sur les fibres de  $X'/S'$. 
Soit  $\alpha$ une famille de  $r$-cycles sur les fibres de  $X/S$  et  
$\alpha' = g^*\alpha$ le changement de base. 
Si  $R$  est un anneau de valuation discrète, alors tout morphisme  
$h : \Spec(R) \to S$  se factorise comme  $g \circ h'$  pour un 
certain morphisme unique  $h' : \Spec(R) \to S'$.
 En effet, le morphisme  $S' \times_S \Spec(R) \to \Spec(R)$  est
 un homéomorphisme universel induisant des bijections sur les corps résiduels,
 et possède donc une section (par exemple parce que  $R$  est un anneau 
semi-normal ; voir Morphismes, section \ref{morphisms-section-seminormalization}). 
Ainsi, la condition que $\alpha$  soit compatible avec les
 spécialisations (c'est-à-dire qu'il s'agisse d'un  $r$-cycle 
relatif) équivaut à la condition que  
$\alpha'$  soit compatible avec les spécialisations. 
\end{proof} 






\section{Cycles relatifs équidimensionnels} 
\label{section-equidimensional} 

\noindent
 Voici la définition. 

\begin{definition} 
\label{definition-equidimensional} 
Soit  $f : X \to S$  un morphisme de schémas. Supposons que  $S$  soit localement noethérien et
 que  $f$  soit localement de type fini. Soit $r \geq 0$  un entier. Nous disons
 qu'un  $r$-cycle relatif  $\alpha$ sur  $X/S$  est {\it équidimensionnel} si le support 
de  $\alpha$  (remarque \ref{remark-supports-family}) est
 contenu dans un sous-ensemble fermé  $W \subset X$  dont la dimension
 relative sur  $S$  est  $\leq r$.
 Le groupe de tous les  $r$-cycles relatifs équidimensionnels sur $X/S$  
est noté  $z_{equi}(X/S, r)$. 
\end{definition} 

\begin{example} 
\label{example-not-equidimensional} 
\begin{reference} 
\cite[Example 3.1.9]{SV} 
\end{reference}
 Il existe des  $r$-cycles relatifs qui ne sont pas équidimensionnels. Plus 
précisément, soit  $k$  un corps et soit  $X = \Spec(k[x, y, t])$
  sur  $S = \Spec(k[x, y])$. Soit $s$  un point de  $S$  et 
notons  $a, b \in \kappa(s)$  les images de $x$  et  $y$.
 Considérons la famille  $\alpha$  de $0$-cycles sur  $X/S$  définie par 
\begin{enumerate} 
\item $\alpha_s = 0$  si  $b = 0$  et sinon 
\item $\alpha_s = [p] - [q]$, où  $p$ (respectivement\ $q$) est le point $\kappa(s)$-rationnel
 de  $\Spec(\kappa(s)[t])$  défini par  $t = a/b$ 
(respectivement\ $t = (a + b^2)/b$). 
\end{enumerate} 
Nous laissons au lecteur le soin de montrer que cela est compatible avec les spécialisations
 ; l'idée est que  $a/b$  et  $(a + b^2)/b = a/b + b$  se spécialisent en un même point 
dans  $\mathbf{P}^1$  sur le corps résiduel de toute valuation  $v$  sur $\kappa(s)$
  avec  $v(b) > 0$. D'autre part, l'adhérence du support de $\alpha$
  contient toute la fibre au-dessus de  $(0, 0)$. 
\end{example} 

\begin{lemma} 
\label{lemma-equidimensional-functoriality} 
Soit  $f : X \to S$  un morphisme de schémas. Supposons que  $S$  soit localement noethérien et
 que  $f$  soit localement de type fini. Soit $r \geq 0$  un entier. Soit  
$\alpha$  un  $r$-cycle relatif sur $X/S$. Si  $\alpha$  est équidimensionnel, 
alors toute restriction, changement de base ou image inverse plate de  $\alpha$  est
 équidimensionnel. 
\end{lemma} 

\begin{proof} La 
démonstration est omise. 
\end{proof} 

\begin{lemma} 
\label{lemma-check-equidimensional} 
Soit  $f : X \to S$  un morphisme de schémas. Supposons  $S$  localement noethérien
 et  $f$  localement de type fini. Soit  $r \geq 0$  un entier. Soit  $\alpha$
  un  $r$-cycle relatif sur  $X/S$. Alors, pour vérifier que  $\alpha$  est équidimensionnel, nous pouvons
 travailler localement pour la topologie de Zariski sur  $X$  et  $S$. 
\end{lemma} 

\begin{proof}
 La condition que  $\alpha$  soit équidimensionnel signifie simplement
 que l'adhérence du support de  $\alpha$  a une dimension relative $\leq r$
  sur  $S$. Comme prendre les adhérences commute avec la restriction aux ouverts,
 le lemme s'ensuit (un détail mineur est omis). 
\end{proof} 

\begin{lemma} 
\label{lemma-equidimensional-fppf-descent} 
Soit  $f : X \to S$  un morphisme de schémas. Supposons  $S$  localement noethérien
 et  $f$  localement de type fini. Soit  $r \geq 0$  un entier. Soit  $\alpha$
  un  $r$-cycle relatif sur  $X/S$. Soit $\{g_i : S_i \to S\}$
  un recouvrement fppf. Alors  $\alpha$  est équidimensionnel si 
et seulement si chaque changement de base  $g_i^*\alpha$  est équidimensionnel. 
\end{lemma} 

\begin{proof} 
Si  $\alpha$  est équidimensionnel, alors chaque  $g_i^*\alpha$  l'est aussi d'après
 le lemme \ref{lemma-equidimensional-functoriality}. Supposons que chaque  $g_i^*\alpha$
  soit équidimensionnel. Notons  $W$  l'adhérence 
de  $\text{Supp}(\alpha)$  dans  $X$. Comme $g_i : S_i \to S$  est universellement ouvert
 (étant plat et localement de présentation finie), il en est de même pour le morphisme  
$f_i : X_i = S_i \times_S X \to X$. Notons  $\alpha_i = g_i^*\alpha$. Nous avons  
$\text{Supp}(\alpha_i) = f_i^{-1}(\text{Supp}(\alpha))$  d'après
 le lemme \ref{lemma-support-family}. 
Puisque  $f_i$  est ouvert, nous voyons que  $W_i = f_i^{-1}(W)$  est l'adhérence 
de  $\text{Supp}(\alpha_i)$. Donc, d'après l'hypothèse, le morphisme  
$W_i \to S_i$  a une dimension relative $\leq r$.
 D'après Morphismes, lemme \ref{morphisms-lemma-dimension-fibre-after-base-change}
 (et le fait que les morphismes $S_i \to S$  sont conjointement surjectifs)
 nous concluons que  $W \to S$  a une dimension relative  $\leq r$. 
\end{proof} 

\begin{lemma} 
\label{lemma-equidimensional-descent-pullbacks} 
Soit  $f : X \to S$  un morphisme de schémas. Supposons  $S$  localement noethérien
 et  $f$  localement de type fini. Soient  $r, e \geq 0$  des entiers. 
Soit  $\alpha$  un  $r$-cycle relatif sur  $X/S$.
 Soit $\{f_i : X_i \to X\}$  une famille conjointement surjective
 de morphismes plats, localement de type fini, et de dimension relative  $e$. 
Alors  $\alpha$  est équidimensionnel si et seulement si chaque image 
inverse plate  $f_i^*\alpha$  est équidimensionnel. 
\end{lemma} 

\begin{proof} La 
démonstration est omise. Indication : comme dans la démonstration du lemme \ref{lemma-equidimensional-fppf-descent},
 on montre que l'image inverse par  $f_i$  de l'adhérence  $W$  du support de  
$\alpha$  est l'adhérence  $W_i$  du support de  $f_i^*\alpha$. Alors, 
$W \to S$  a une dimension relative  $\leq r$  si  $W_i \to S$
  a une dimension relative  $\leq r + e$  pour tout  $i$. 
\end{proof} 





\section{Pondérations et zéro-cycles relatifs} 
\label{section-weightings} 

\noindent
 Dans cette section, nous montrons que pour un morphisme quasi-fini, les 
zéro-cycles relatifs sont étroitement liés aux pondérations, telles qu'elles sont définies dans 
Compléments sur les morphismes, section \ref{more-morphisms-section-weightings}. 

\medskip\noindent
 Soit  $S$  un schéma localement noethérien. Soit  $f : X \to S$  un morphisme 
localement quasi-fini de schémas. Alors nous avons  $z(X/S, 0) = z_{equi}(X/S, 0)$
  et  $z(X/S, r) = 0$ pour  $r > 0$. Étant donné  $\alpha \in z(X/S, 0)$ , définissons
 une application 
$$
w_\alpha : X \longrightarrow \mathbf{Z},\quad
x \mapsto \alpha(x) [\kappa(x) : \kappa(s)]_i \quad\text{où }s = f(x)
$$
 Ici,  $\alpha(x)$  désigne le coefficient de  $x$  dans le $0$-cycle  
$\alpha_s$  sur la fibre  $X_s$  et  $[K : k]_i$  désigne le degré inséparable
 d'une extension finie de corps. 

\begin{lemma} 
\label{lemma-weightings-pre} 
Dans la situation ci-dessus, si  $g : S' \to S$  est un morphisme de schémas 
localement noethériens, alors  $w_{g^*\alpha} = w_\alpha \circ g'$  où  
$g' : X' \to X$  est la projection $X' = S' \times_S X \to X$. 
\end{lemma} 

\begin{proof} 
Soit  $x' \in X'$  dont les images sont  $s', s, x$  dans  $S', S, X$.
 Alors le coefficient de  $[x']$  dans
 le changement de base de  $[x]$  par $\kappa(s')/\kappa(s)$  est la longueur de
 l'anneau local  $(\kappa(s') \otimes_{\kappa(s)} \kappa(x))_\mathfrak q$.
 Ici  $\mathfrak q$  est l'idéal premier correspondant à  $x'$. 
Ainsi, la compatibilité avec le changement de base résulte de l'égalité 
$$
[\kappa(x) : \kappa(s)]_i =
\text{longueur}((\kappa(s') \otimes_{\kappa(s)} \kappa(x))_\mathfrak q)
[\kappa(x') : \kappa(s')]_i
$$
 Soit l'extension  $k/\kappa(s')$  une clôture algébrique. Choisissons un idéal premier  
$\mathfrak p \subset k \otimes_{\kappa(s)} \kappa(x)$ au-dessus
 de  $\mathfrak q$. Supposons que nous puissions montrer que 
$$
[\kappa(x) : \kappa(s)]_i =
\text{longueur}((k \otimes_{\kappa(s)} \kappa(x))_\mathfrak p)
\quad\text{et}\quad
[\kappa(x') : \kappa(s')]_i =
\text{longueur}((k \otimes_{\kappa(s')} \kappa(x'))_\mathfrak p)
$$
 Le résultat s'ensuit car 
$$
\text{longueur}((\kappa(s') \otimes_{\kappa(s)} \kappa(x))_\mathfrak q)
\text{longueur}((k \otimes_{\kappa(s')} \kappa(x'))_\mathfrak p)
=
\text{longueur}((k \otimes_{\kappa(s)} \kappa(x))_\mathfrak p)
$$
 d'après Algèbre, lemme \ref{algebra-lemma-pullback-module}, et la platitude de  
$\kappa(s') \otimes_{\kappa(s)} \kappa(x) \to k \otimes_{\kappa(s)} \kappa(x)$.
 Pour montrer les deux égalités, il suffit de prouver la première. 
Soit  $\kappa(x)/\kappa/\kappa(s)$  le sous-corps construit dans
 Corps, lemme \ref{fields-lemma-separable-first}. On obtient alors 
$$
k \otimes_{\kappa(s)} \kappa(x) =
\prod\nolimits_{\sigma : \kappa \to k}
k \otimes_{\sigma, \kappa} \kappa(x)
$$
 et chacun des facteurs est local de degré  
$[\kappa(x) : \kappa] = [\kappa(x) : \kappa(s)]_i$
 , comme voulu. 
\end{proof} 

\begin{lemma} 
\label{lemma-weightings} 
Soit  $S$  un schéma localement noethérien. Soit  $f : X \to S$  un morphisme
 de schémas localement quasi-fini. 
\begin{enumerate} 
\item Pour  $\alpha \in z(X/S, 0)$ , l'application  $w_\alpha : X \to \mathbf{Z}$
  construite ci-dessus est une pondération. 
\item Si  $X$  est quasi-compact, alors étant donnée une pondération 
$w : X \to \mathbf{Z}$ , il existe un entier  $n > 0$  tel 
que  $nw = w_\alpha$  pour un certain  $\alpha \in z(X/S, 0)$. 
\item L'entier  $n$  de (2) peut être choisi égal à une puissance du 
nombre premier  $p$  si  $S$  est un schéma sur  $\mathbf{F}_p$. 
\end{enumerate} 
\end{lemma} 

\begin{proof} 
Soit  $\alpha \in z(X/S, 0)$  et choisissons un diagramme 
$$
\xymatrix{
X \ar[d]_f & U \ar[l]^h \ar[d]^\pi \\
Y & V \ar[l]_g
}
$$
 comme dans Compléments sur les morphismes, définition \ref{more-morphisms-definition-weighting}. 
Notons  $\beta \in z(U/V, 0)$  la restriction du changement de base  $g^*\alpha$.
 Par la compatibilité avec le changement de base (lemme \ref{lemma-weightings-pre}), 
nous avons  $w_\beta = w_\alpha \circ h$  et il suffit de montrer que  
$\int_\pi w_\beta$  est localement constant sur  $V$. Ensuite, notons que 
\begin{align*} 
\left( \int_\pi w_\beta \right)(v) 
& = 
\sum\nolimits_{u \in U, \pi(u) = v} 
\beta(u) [\kappa(u) : \kappa(v)]_i [\kappa(u) : \kappa(v)]_s \\
& = 
\sum\nolimits_{u \in U, \pi(u) = v} \beta(u)[\kappa(u) : \kappa(v)] 
\end{align*} 
Cette dernière expression est le coefficient de  $v$  dans  $\pi_*\beta \in z(V/V, 0)$. 
D'après le lemme \ref{lemma-relative-cycle-smooth}, cette fonction est localement 
constante sur $V$. 

\medskip\noindent
 Réciproquement, soit  $w : X \to S$  une pondération et  $X$  quasi-compact.
 Choisissez un entier  $n$ suffisamment divisible. Soit  $\alpha$  la famille 
de  $0$-cycles sur les fibres de  $X/S$  telle que pour $s \in S$  nous avons 
$$
\alpha_s =
\sum\nolimits_{f(x) = s} \frac{n w(x)}{[\kappa(x) : \kappa(s)]_i} [x]
$$
 comme un zéro-cycle sur  $X_s$.
 Cela a un sens puisque les fibres de $f$  sont universellement bornées
 (Morphismes, lemme 
\ref{morphisms-lemma-locally-quasi-finite-qc-source-universally-bounded}), donc 
nous pouvons trouver  $n$  de sorte que le membre de droite soit un entier pour 
tout  $s \in S$. L'énoncé final du lemme en découle également, à condition que nous montrions
 que  $\alpha$  est un  $0$-cycle relatif. Pour ce faire, il faut 
montrer que  $\alpha$  est compatible avec les spécialisations le long des spectres
 d'anneaux de valuation discrète. Notons que ni la construction de $w_\alpha$
  ni la démonstration du lemme \ref{lemma-weightings-pre} ne supposent que  $\alpha$
 soit un cycle relatif. Par conséquent, nous avons  $w_\alpha = nw$  et cela reste vrai 
après un changement de base. De même, le changement de base d'une pondération est une
 pondération ; voir 
Compléments sur les morphismes, lemme \ref{more-morphisms-lemma-weighting-base-change}. 
Ainsi nous revenons au problème étudié dans le paragraphe suivant. 

\medskip\noindent
 Supposons que  $S$  soit le spectre d'un anneau de valuation discrète avec un point 
générique  $\eta$  et un point fermé  $0$. Soit  $w : X \to S$ une pondération, avec  
$X$  quasi-fini sur  $S$. Soit  $\alpha$  la famille de  $0$-cycles sur les
 fibres de  $X/S$  construite dans le paragraphe précédent (pour un  $n$ 
approprié). Nous devons montrer que $sp_{X/S}(\alpha_\eta) = \alpha_0$. 
Soit  $\beta \in z(X/S, 0)$  le  $0$-cycle relatif sur  $X/S$
 avec  $\beta_\eta = \alpha_\eta$  et  $\beta_0 = sp_{X/S}(\alpha_\eta)$. 
Alors  $w' = w_\beta - nw : X \to \mathbf{Z}$
  est une pondération (en utilisant le résultat ci-dessus) et nulle aux points de  $X$
  qui se projettent sur  $\eta$. Il résulte de Compléments sur les morphismes,
 lemme \ref{more-morphisms-lemma-weighting-specialization} 
que  $w' = 0$. Cela signifie  $w_\beta = nw$, donc  $\alpha = \beta$  comme souhaité. 
\end{proof} 

\begin{example} 
\label{example-n-must-be-positive-sometimes} 
Soit  $p$  un nombre premier. Soit  $k'/k$  une extension purement inséparable
 de corps de degré  $p^f$. Posons  $X = \Spec(k')$  et  $S = \Spec(k)$. 
L'application  $w : X \to \mathbf{Z}$  envoyant le point unique  $x$  de  $X$ sur  $1$
  est une pondération de  $X \to S$. D'autre part, 
le groupe $z(X/S, 0)$  est librement engendré par le cycle  $\alpha = [x]$.
 Ainsi, nous voyons que le plus petit entier  $n$  qui convient dans 
le lemme \ref{lemma-weightings} est  $n = p^f$  dans ce cas. 
\end{example} 

\begin{example} 
\label{example-Merkurjev} La
 discussion ci-dessus peut être utilisée pour « expliquer » un exemple dû 
à A.S.~Merkurjev \cite[Example 3.5.10]{SV}. Soit  $p$  un nombre premier 
et soit  $k = \mathbf{F}_p(a, b)$  l'extension purement transcendante 
de  $\mathbf{F}_p$ engendrée par  $a$  et  $b$. Posons 
$$
A = k[x, y, z]/(a x^p + b y^p - z^p)
$$
 L'anneau ainsi obtenu est un anneau intègre normal de type fini sur  $k$. Le spectre  
$S = \Spec(A)$ est régulier, à l'exception d'un point singulier unique  $s$
  correspondant à l'idéal maximal  $(x, y, z)$  de  $A$, de corps résiduel  $k$. 
Soient  $k' = k(a^{1/p}, b^{1/p})$  et  $B = k'[x, y]$. Considérons l'homomorphisme  
$A \to B$ induit par l'inclusion  $k \to k'$ et envoyant  $x$  sur  $x$, $y$  sur  $y$,
 et  $z$  sur  $a^{1/p}x + b^{1/p}y$. Posons $X = \Spec(B)$  et considérons le morphisme  
$f : X \to S$  correspondant au morphisme d'anneaux  $A \to B$.
 La pondération  $w : X \to \mathbf{Z}$  de  $f$
  donnée par Compléments sur les morphismes, lemme 
\ref{more-morphisms-lemma-weighting-quasi-finite-Noetherian}, est 
constante de valeur  $1$, puisque le corps des fractions de  $B$  est une extension purement inséparable de 
celui de  $A$. 
Posons  $V = f^{-1}(U) \subset X$. Par le critère de platitude par miracle, le morphisme  
$V \to U$ est plat de degré  $p$. Il s'ensuit que le cycle  
$\alpha = [V/V/U]_0 \in z(V/U, 0)$  satisfait $w_\alpha = pw|_V$.
 Cependant, puisque le point unique  $x$  de  $X$ situé au-dessus de  $s$
  a pour corps résiduel  $k'$  de degré $p^2$  sur  $k$, nous voyons 
qu'il ne peut y avoir d'élément $\beta$  de  $z(X/S, 0)$  se restreignant à  $\alpha$.
 En effet, la pondération $pw - w_\beta$  serait nulle d'après Compléments sur les morphismes,
 lemme \ref{more-morphisms-lemma-weighting-specialization}, ce qui 
impliquerait (d'après la formule dans la démonstration du lemme ci-dessus) que 
$$
\beta_s = \frac{pw(x)}{[\kappa(x) : \kappa(s)]_i} [x] =
\frac{1}{p}[x]
$$
 dans  $Z_0(X_s)$  et c'est impossible\footnote{Dans  \cite{SV}  il existe
 un élément du groupe des cycles relatifs  $Cycl(X/S, 0)$
  correspondant à $pw$  et ce cycle se spécialise en un cycle sur
  $X_s$  avec des coefficients non entiers.}. Bien sûr,
 le cycle  $p\alpha$  est la restriction d'un élément unique de  
$z(X/S, 0)$. 
\end{example} 









\section{Cycles relatifs effectifs} 
\label{section-effective} 


\noindent
 Voici la définition. 

\begin{definition} 
\label{definition-effective} 
Soit  $f : X \to S$  un morphisme de schémas. Supposons que  $S$  soit localement noethérien et
 que  $f$  soit localement de type fini. Soit $r \geq 0$  un entier. On dit 
qu'un  $r$-cycle relatif  $\alpha$  sur $X/S$  est {\it effectif} si  $\alpha_s$ est un
 cycle effectif 
(Homologie de Chow, définition \ref{chow-definition-effective-cycle}) 
pour tout $s \in S$. Le monoïde de tous les  $r$-cycles relatifs effectifs
 sur $X/S$  est noté  $z^{eff}(X/S, r)$. 
\end{definition} 

\noindent
 Plus loin, nous montrerons qu'un cycle relatif effectif est équidimensionnel ; 
voir le lemme \ref{lemma-effective-equidimensional}. 

\begin{lemma} 
\label{lemma-effective-functoriality} 
Soit  $f : X \to S$  un morphisme de schémas. Supposons que  $S$  soit localement noethérien et
 que  $f$  soit localement de type fini. Soit $r \geq 0$  un entier. Soit  
$\alpha$  un  $r$-cycle relatif sur $X/S$. Si  $\alpha$  est effectif, alors 
l'effectivité est préservée par toute restriction, tout changement de base, toute image inverse plate et toute image directe propre
 de  $\alpha$. 
\end{lemma} 

\begin{proof} La 
démonstration est omise. 
\end{proof} 

\begin{lemma} 
\label{lemma-check-effective} 
Soit  $f : X \to S$  un morphisme de schémas. Supposons que  $S$  soit localement noethérien et
 que  $f$  soit localement de type fini. Soit $r \geq 0$  un entier. Soit  $\alpha$
  un  $r$-cycle relatif sur $X/S$. Alors pour vérifier que  $\alpha$  est effectif, nous pouvons
 travailler localement pour la topologie de Zariski sur  $X$  et $S$. 
\end{lemma} 

\begin{proof} La 
démonstration est omise. 
\end{proof} 

\begin{lemma} 
\label{lemma-effective-descent} 
Soit $f : X \to S$ un morphisme de schémas. Supposons $S$ localement noethérien
 et $f$ localement de type fini. Soit $r \geq 0$ un entier. Soit $\alpha$
 un $r$-cycle relatif sur $X/S$. Soit $g : S' \to S$ un morphisme surjectif. 
Alors $\alpha$ est effectif si et seulement si le changement de base $g^*\alpha$
 est effectif. 
\end{lemma} 

\begin{proof} La 
démonstration est omise. 
\end{proof} 

\begin{lemma} 
\label{lemma-effective-descent-pullbacks} 
Soit $f : X \to S$ un morphisme de schémas. Supposons $S$ localement noethérien
 et $f$ localement de type fini. Soient $r, e \geq 0$ des entiers. 
Soit $\alpha$ un $r$-cycle relatif sur $X/S$.
 Soit $\{f_i : X_i \to X\}$ une famille conjointement surjective
 de morphismes plats, localement de type fini et de dimension relative $e$.
 Alors $\alpha$ est effectif si et seulement si chaque
 cycle $f_i^*\alpha$ obtenu par image inverse plate est effectif. 
\end{lemma} 

\begin{proof} La 
démonstration est omise. 
\end{proof} 

\begin{lemma} 
\label{lemma-effective-support-closed} 
Soit  $f : X \to S$  un morphisme de schémas. Supposons que  $S$  soit localement noethérien
 et  $f$  localement de type fini. Soient  $r, e \geq 0$ des entiers. 
Soit  $\alpha$  un  $r$-cycle relatif sur  $X/S$.
 Si $\alpha$  est effectif, alors  $\text{Supp}(\alpha)$  est
 fermé dans  $X$. 
\end{lemma} 

\begin{proof} 
Soit  $g : S' \to S$  l'inclusion d'une composante irréductible vue 
comme un sous-schéma fermé intègre. D'après les 
lemmes  \ref{lemma-effective-functoriality}  et  \ref{lemma-support-family} ,
 il suffit de montrer que le support du changement de base  
$g^*\alpha$  est fermé dans $S' \times_S S$. 
Ainsi nous pouvons supposer que  $S$  est un schéma intègre de point générique  
$\eta$. Nous montrerons que  $\text{Supp}(\alpha)$  est l'adhérence 
de  $\text{Supp}(\alpha_\eta)$. Pour ce faire, choisissons un point  $s \in S$. 
Nous pouvons trouver un morphisme  $g : S' \to S$  où  $S'$  est le spectre 
d'un anneau de valuation discrète envoyant le point générique  $\eta' \in S'$  vers $\eta$
  et le point fermé  $0 \in S'$  vers  $s$, voir
 Propriétés, lemme \ref{properties-lemma-locally-Noetherian-specialization-dvr}. 
Il suffit alors de prouver que le support de  $g^*\alpha$
  est égal à l'adhérence de  $\text{Supp}((g^\alpha)_{\eta'})$.
 Cela nous ramène au cas discuté dans le paragraphe suivant. 

\medskip\noindent
 Ici  $S$  est le spectre d'un anneau de valuation discrète avec point 
générique  $\eta$  et point fermé  $0$. Nous devons montrer que 
$\text{Supp}(\alpha)$  est l'adhérence de  $\text{Supp}(\alpha_\eta)$. 
Puisque  $\alpha$  est effectif, nous pouvons écrire  $\alpha_\eta = \sum n_i[Z_i]$
  avec  $n_i > 0$  et les  $Z_i \subset X_\eta$ sont des sous-schémas fermés intègres de dimension  $r$. 
Puisque  $\alpha_0 = sp_{X/S}(\alpha_\eta)$  nous savons que 
$\alpha_0 = \sum n_i [\overline{Z}_{i, 0}]_r$  où  $\overline{Z}_i$
  est l'adhérence de  $Z_i$. D'après Variétés, lemme 
\ref{varieties-lemma-dominate-valuation-ring-dimension-fibres}, nous
 voyons que  $\overline{Z}_{i, 0}$  est équidimensionnel de dimension $r$. 
Puisque  $n_i > 0$ , nous en concluons que  $\text{Supp}(\alpha_0)$
  est égal à l'union des  $\overline{Z}_{i, 0}$, c'est-à-dire
 à la fibre au-dessus de  $0$  de $\bigcup \overline{Z}_i$, laquelle
 est à son tour l'adhérence de  $\bigcup Z_i$, comme voulu. 
\end{proof} 

\begin{lemma} 
\label{lemma-effective-equidimensional} 
Soit  $f : X \to S$  un morphisme de schémas. Supposons  $S$  localement noethérien
 et  $f$  localement de type fini. Soient  $r, e \geq 0$  des entiers. 
Soit  $\alpha$  un  $r$-cycle relatif sur  $X/S$.
 Si $\alpha$  est effectif, alors  $\alpha$  est équidimensionnel. 
\end{lemma} 

\begin{proof} 
Supposons que  $\alpha$  soit effectif. D'après le lemme \ref{lemma-effective-support-closed},
 le support $\text{Supp}(\alpha)$  est fermé dans  $X$. Ainsi,  
$\alpha$  est équidimensionnel car les fibres de  $\text{Supp}(\alpha) \to S$
  sont les supports des cycles  $\alpha_s$  et ont donc dimension  $r$. 
\end{proof} 

\begin{remark} 
\label{remark-representable} 
Soit  $f : X \to S$  un morphisme de schémas avec  $S$  localement noethérien
 et  $f$  localement de type fini. On peut se demander si le foncteur contravariant 
$$
\begin{matrix}
\text{schémas }S'\text{ localement} \\
\text{de type fini sur }S
\end{matrix}
\longrightarrow
z^{eff}(X'/S', r)\text{ où }X' = S' \times_S X
$$
 est représentable. Comme  $z(X'/S', r) = z(X'_{red}/S'_{red}, r)$
 , cela ne peut pas être vrai (nous laissons au lecteur le soin de donner un véritable
 contre-exemple). Une meilleure question serait de savoir si nous pouvons 
trouver une sous-catégorie du membre de gauche sur laquelle le foncteur est
 représentable. Le lemme  \ref{lemma-seminormalize-base}  suggère que 
nous devrions au moins nous restreindre à la catégorie des schémas semi-normaux sur $S$. 

\medskip\noindent
 Si  $S/\Spec(\mathbf{Q})$  est Nagata et  $f$  est un morphisme projectif, alors
 il s'avère que  $S' \mapsto z^{eff}(X'/S', r)$  est représentable sur la 
catégorie dont les objets $S'$  sont semi-normaux. Essentiellement, c'est le contenu de 
\cite[Theorem 3.21]{KRC}. 

\medskip\noindent
 Si  $S$  possède des points de caractéristique positive, cela ne fonctionne plus même
 si nous remplaçons la semi-normalité par la normalité faible ; un schéma localement 
noethérien  $T$  est faiblement normal si tout homéomorphisme universel birationnel  
$T' \to T$  possède une section. Un exemple consiste à considérer les  $0$-cycles de degré  
$2$  sur  $X = \mathbf{A}^2_k$ au-dessus de  $S = \Spec(k)$  où  $k$  est un corps 
de caractéristique $2$. Plus précisément, au-dessus de  $W = X \times_S X$  nous avons
 un  $0$-cycle relatif canonique  $\alpha \in z^{eff}(X_W/W, 0)$ : pour  
$w = (x_1, x_2) \in W = X^2$  nous avons le cycle $\alpha_w = [x_1] + [x_2]$.
 Ce cycle est invariant sous l'involution  $\sigma : W \to W$  échangeant
 les facteurs. Puisque  $W$  est lisse (donc normal, donc faiblement normal), si 
$z(-/-, r)$  était représentable par  $M$  sur la catégorie des schémas faiblement
 normaux de type fini sur  $k$ , nous obtiendrions un morphisme invariant par  $\sigma$, 
de  $W$  vers  $M$. Cela définirait à son tour un morphisme du 
schéma quotient  $\text{Sym}^2_S(X) = W/\langle \sigma \rangle$  vers  $M$.
 Puisque $\text{Sym}^2_S(X)$  est normal, nous obtiendrions, par la propriété universelle de l'espace de modules $M$
 , un  $0$-cycle relatif  $\beta$  sur  
$X \times_S \text{Sym}^2_S(X) / \text{Sym}^2_S(X)$
  dont l'image inverse sur  $W$  est  $\alpha$. Cependant, il n'existe aucun cycle $\beta$  
de cette sorte. En effet, en écrivant  $X = \Spec(k[u, v])$, le schéma  
$\text{Sym}^2_S(X)$ est le spectre de 
$$
k[u_1 + u_2, u_1u_2, v_1 + v_2, v_1v_2, u_1v_1 + u_2v_2]
\subset
k[u_1, u_2, v_1, v_2]
$$
 L'image de la diagonale  $u_1 = u_2, v_1 = v_2$  dans  $\text{Sym}^2_S(X)$
  est le sous-schéma fermé $V = \Spec(k[u_1^2, v_1^2])$; ici nous utilisons le 
fait que la caractéristique de  $k$ est  $2$. En regardant le 
point générique  $\eta$  de  $V$, le cycle  $\beta_\eta$  serait un 
zéro-cycle de degré  $2$  sur  $\mathbf{A}^2_{k(u_1^2, v_1^2)}$
 dont l'image inverse sur  $\mathbf{A}^2_{k(u_1, u_2)}$  serait  
$2[\text{le point de coordonnées} (u_1, v_2)]$.
 Cela est clairement impossible. 

\medskip\noindent
 La discussion ci-dessus ne contredit pas  \cite[Theorem 4.13]{KRC}  car la variété
 de Chow de ce théorème ne représente un foncteur qu'au sens grossier (en
 fait deux foncteurs distincts, dont un seul coïncide avec le nôtre pour  $X$
  projectif comme on peut le voir avec un peu de travail). De même, dans  \cite[Section 4.4]{SV}
  il est montré que pour  $X/S$  projectif, la faisceautisation pour la topologie  $h$ du préfaisceau 
$S' \mapsto z^{eff}(S' \times_S X/S', r)$  est égale à la faisceautisation pour la topologie  $h$ d'un foncteur
 représentable. 
\end{remark} 

\begin{remark} 
\label{remark-equidimensional-over-geometrically-unibranch} 
Soit  $f : X \to S$  un morphisme de schémas. Soit  $r \geq 0$. Soit  $Z \subset X$
  un sous-schéma fermé. Supposons que 
\begin{enumerate} 
\item $S$  soit noethérien et géométriquement unibranche, que 
\item $f$  soit de type fini, et que 
\item $Z \to S$  ait une dimension relative  $\leq r$. 
\end{enumerate} 
Alors, pour tout entier suffisamment divisible $n \geq 1$, il 
existe un unique $r$-cycle relatif effectif $\alpha$ sur $X/S$ tel que 
$\alpha_\eta = n[Z_\eta]_r$ pour tout point générique $\eta$  de  $S$.
 Il s'agit d'une reformulation de  \cite[Theorem 3.4.2]{SV}. 
Si nous avons besoin de ce résultat, nous le préciserons et le prouverons ici. 
\end{remark} 











\section{Cycles relatifs propres} 
\label{section-proper} 

\noindent
 Dans notre contexte, la définition suivante est probablement la bonne. 

\begin{definition} 
\label{definition-proper} 
Soit  $f : X \to S$  un morphisme de schémas. Supposons que  $S$  soit localement noethérien et
 que  $f$  soit localement de type fini. Soit $r \geq 0$  un entier. On dit 
qu'un  $r$-cycle relatif  $\alpha$  sur $X/S$  est un  {\it  cycle relatif propre} 
 si le support de  $\alpha$  (remarque \ref{remark-supports-family}) est
 contenu dans un sous-ensemble fermé  $W \subset X$  propre sur  $S$
  (Cohomologie des schémas, définition \ref{coherent-definition-proper-over-base}).
 Le groupe de tous les  $r$-cycles relatifs propres sur  $X/S$  
est noté  $c(X/S, r)$. 
\end{definition} 

\noindent
 D'après Cohomologie des schémas, lemme 
\ref{coherent-lemma-closed-closed-proper-over-base}, cela 
signifie simplement que l'adhérence du support est propre sur la base. Pour voir que
 ces cycles forment un groupe, on utilise
 Cohomologie des schémas, lemme \ref{coherent-lemma-union-closed-proper-over-base}. 

\begin{lemma} 
\label{lemma-proper-functoriality} 
Soit  $f : X \to S$  un morphisme de schémas. Supposons que  $S$  soit localement noethérien et
 que  $f$  soit localement de type fini. Soit $r \geq 0$  un entier. Soit  
$\alpha$  un  $r$-cycle relatif sur $X/S$. Si  $\alpha$  est propre, 
alors tout changement de base de $\alpha$  est propre. 
\end{lemma} 

\begin{proof} La 
démonstration est omise. 
\end{proof} 

\begin{lemma} 
\label{lemma-proper-h-descent} 
Soit  $f : X \to S$  un morphisme de schémas. Supposons que  $S$  soit localement noethérien et
 que  $f$  soit localement de type fini. Soit $r \geq 0$  un entier. Soit  $\alpha$
  un  $r$-cycle relatif sur $X/S$. Soit  $\{g_i : S_i \to S\}$
  un recouvrement h. Alors  $\alpha$  est propre 
si et seulement si chaque changement de base  $g_i^*\alpha$  est propre. 
\end{lemma} 

\begin{proof} 
Si  $\alpha$  est propre, alors chaque  $g_i^*\alpha$  l'est aussi d'après
 le lemme \ref{lemma-proper-functoriality}. Supposons que chaque  $g_i^*\alpha$
  soit propre. Pour prouver que $\alpha$  est propre, il suffit de travailler 
localement sur des ouverts affines sur  $S$. Nous pouvons donc supposer que  $S$
  est affine. Nous pouvons alors affiner notre recouvrement  $\{S_i \to S\}$
  par une famille $\{T_j \to S\}$  où  $g : T \to S$  est un morphisme propre 
et surjectif et  $T = \bigcup T_j$ est un recouvrement ouvert.
 Il s'ensuit que  $\beta = g^*\alpha$  est propre sur $Y = T \times_S X$
  au-dessus de  $T$. D'après le lemme \ref{lemma-support-family}, le support de
  $\beta$  est l'image inverse du support de $\alpha$  par le 
morphisme  $f : Y \to X$. Ainsi, l'adhérence  $W \subset Y$
  de $f^{-1}\text{Supp}(\alpha)$  est propre sur  $T$. Puisque le 
morphisme  $T \to S$  est propre, il s'ensuit que  $W$  est propre sur  $S$.
 D'après Cohomologie des schémas, lemme 
\ref{coherent-lemma-functoriality-closed-proper-over-base},
 l'image  $f(W) \subset X$  est un sous-ensemble fermé propre sur  $S$. 
Puisque  $f(W)$  contient $\text{Supp}(\alpha)$ , nous concluons que  $\alpha$
  est propre. 
\end{proof} 



\section{Cycles relatifs propres et équidimensionnels} 
\label{section-proper-equidimensional} 

\noindent
 Soit  $f : X \to S$  un morphisme de schémas. Supposons que  $S$  soit localement noethérien et
 que  $f$  soit localement de type fini. Soit $r \geq 0$  un entier. On dit 
qu'un  $r$-cycle relatif  $\alpha$  sur $X/S$  est un {\it cycle relatif propre et 
équidimensionnel} si  $\alpha$  est à la fois équidimensionnel
 (définition \ref{definition-equidimensional})
 et propre (définition \ref{definition-proper}).
 Le groupe de tous les  $r$-cycles relatifs propres et équidimensionnels sur  $X/S$  
est noté  $c_{equi}(X/S, r)$. 

\medskip\noindent
 De même, nous disons qu'un  $r$-cycle relatif  $\alpha$  sur $X/S$  est un
 {\it cycle relatif propre et effectif} si  $\alpha$  est à la fois effectif
 (définition \ref{definition-effective}) et propre 
(définition \ref{definition-proper}).
 Le monoïde de tous les $r$-cycles relatifs propres et effectifs sur  $X/S$  
est noté $c^{eff}(X/S, r)$.
 Ces cycles sont équidimensionnels d'après le
 lemme \ref{lemma-effective-equidimensional}. 

\medskip\noindent
 Ainsi, nous avons le diagramme suivant des applications d'inclusion 
$$
\xymatrix{
c^{eff}(X/S, r) \ar[r] \ar[d] &
c_{equi}(X/S, r) \ar[r] \ar[d] &
c(X/S, r) \ar[d] \\
z^{eff}(X/S, r) \ar[r] &
z_{equi}(X/S, r) \ar[r] &
z(X/S, r)
}
$$



\section{Action sur les cycles} 
\label{section-action} 

\noindent
 Soit  $S$  un schéma localement noethérien, universellement caténaire, muni 
d'une fonction de dimension  $\delta$ ; voir
 Homologie de Chow, section \ref{chow-section-setup}. 
Soit  $X \to Y$  un morphisme de schémas sur  $S$, tous deux localement de type
 fini sur  $S$. Soit  $r \geq 0$. Enfin, soit $\alpha$  une famille de  
$r$-cycles sur les fibres de  $X/Y$. Pour  $e \in \mathbf{Z}$
 , nous allons construire une opération 
$$
\alpha \cap - : Z_e(Y) \longrightarrow Z_{r + e}(X)
$$
 Concrètement, étant donné  $\beta \in Z_e(Y)$, écrivons  $\beta = \sum n_i[Z_i]$
  où  $Z_i \subset Y$  est un sous-schéma fermé intègre de  
$\delta$-dimension  $e$  et la 
famille $Z_i$  est localement finie dans le schéma  $Y$. 
Soit  $y_i \in Z_i$ le point générique. Écrivons  
$\alpha_{y_i} = \sum m_{ij} [V_{ij}]$. Ainsi  $V_{ij} \subset X_{y_i}$
  est un sous-schéma fermé intègre de dimension  $r$  et la famille 
$V_{ij}$ est localement finie dans le schéma $X_{y_i}$. 
Nous posons alors 
$$
\alpha \cap \beta = \sum n_i m_{ij} [\overline{V}_{ij}]
\quad\in\quad
Z_{r + e}(X)
$$
 Ici, $\overline{V}_{ij} \subset X$ est l'image schématique 
du morphisme $V_{ij} \to X_{y_i} \to X$ ou, de manière équivalente, 
$\overline{V}_{ij} \subset X$ est un sous-schéma fermé 
intègre qui domine $Z_i \subset Y$ et dont la fibre générique est $V_{ij}$.
 On en déduit aisément que $\dim_\delta(\overline{V}_{ij}) = r + e$
 et que la famille de sous-schémas fermés 
$\overline{V}_{ij} \subset X$ est localement finie (nous omettons les vérifications). 
Ainsi, $\alpha \cap \beta$ est bien un élément de $Z_{r + e}(X)$. 

\begin{lemma} 
\label{lemma-action-bilinear} La 
construction ci-dessus est bilinéaire, c'est-à-dire que 
$(\alpha_1 + \alpha_2) \cap \beta = \alpha_1 \cap \beta +
\alpha_2 \cap \beta$ et $\alpha \cap (\beta_1 + \beta_2) =
\alpha \cap \beta_1 + \alpha \cap \beta_2$. 
\end{lemma} 

\begin{proof} 
Démonstration omise. 
\end{proof} 

\begin{lemma} 
\label{lemma-action-opens} 
Si $U \subset X$ et $V \subset Y$ sont ouverts et si $f(U) \subset V$, alors 
$(\alpha \cap \beta)|_U$ est égal à $\alpha|_U \cap \beta|_V$. 
\end{lemma} 

\begin{proof} 
Cela résulte immédiatement de la description explicite de $\alpha \cap \beta$
 donnée ci-dessus. 
\end{proof} 

\begin{lemma} 
\label{lemma-action-base-change} La 
formation de $\alpha \cap \beta$ est compatible avec le changement de base plat et avec 
l'image inverse plate (voir la démonstration pour plus de précisions). 
\end{lemma} 

\begin{proof} 
Soient $(S, \delta)$, $(S', \delta')$, $g : S' \to S$ et $c \in \mathbf{Z}$
 comme dans Homologie de Chow, situation \ref{chow-situation-setup-base-change}. 
Soit $X \to Y$ un morphisme de schémas localement de type fini sur $S$.
 Notons $X' \to Y'$ le changement de base de $X \to Y$ par $g$.
 Soit $\alpha$ une famille de $r$-cycles sur les fibres de $X/Y$.
 Soit $\beta \in Z_e(Y)$. Notons $\alpha'$ le 
changement de base de $\alpha$ par $Y' \to Y$. Notons 
$\beta' = g^*\beta \in Z_{e + c}(Y')$ l'image inverse de $\beta$ par $g$; voir
 Homologie de Chow, section \ref{chow-section-change-base}. La 
compatibilité avec le changement de base signifie que 
$\alpha' \cap \beta'$ est le changement de base de $\alpha \cap \beta$. 

\medskip\noindent
 Démonstration de la compatibilité au changement 
de base. Puisque nous démontrons une égalité de cycles sur $X'$, nous pouvons travailler localement
 sur $Y$; voir le lemme \ref{lemma-action-opens}. Nous pouvons donc supposer $Y$
 affine. En particulier, $\beta$ est une combinaison linéaire finie de cycles 
premiers. Puisque $- \cap -$ est linéaire en la seconde variable
 (lemme \ref{lemma-action-bilinear}), il suffit de 
démontrer l'égalité lorsque $\beta = [Z]$ pour un sous-schéma fermé intègre 
$Z \subset Y$ de dimension $\delta$ égale à $e$. 

\medskip\noindent
 Soit $y \in Z$ le point générique. Écrivons $\alpha_y = \sum m_j [V_j]$.
 Soit $\overline{V}_j$ l'adhérence de $V_j$ dans $X$. Alors 
$$
\alpha \cap \beta = \sum m_j[\overline{V}_j]
$$
 Le changement de base de $\beta$ est $\beta' = \sum [Z \times_S S']_{e + c}$, considéré
 comme un cycle sur $Y' = Y \times_S S'$. Soient $Z'_a \subset Z \times_S S'$
 les composantes irréductibles; notons $y'_a \in Z'_a$ leurs points 
génériques et $n_a$ la multiplicité de $Z'_a$ dans $Z \times_S S'$. 
Nous avons 
$$
\beta' = \sum [Z \times_S S']_{e + c} = \sum n_a[Z'_a]
$$
 Nous avons $\alpha'_{y'_a} = \sum m_j [V_{j, \kappa(y'_a)}]_r$
 car $\alpha'$ est le changement de base de $\alpha$ par $Y' \to Y$. 
Soient $V'_{jab} \subset V_{j, \kappa(y'_a)}$ les composantes irréductibles,
 et notons $m_{jab}$ la multiplicité de $V'_{jab}$
 dans $V_{j, \kappa(y'_a)}$. Nous avons 
$$
\alpha'_{y'_a} = \sum m_j [V_{j, \kappa(y'_a)}]_r =
\sum m_j m_{jab} [V'_{jab}]
$$
 Ainsi, 
$$
\alpha' \cap \beta' = \sum n_a m_j m_{jab} [\overline{V}'_{jab}]
$$
 où $\overline{V}'_{jab}$ est l'adhérence de $V'_{jab}$ dans $X'$.
 Pour démontrer l'égalité voulue, il suffit donc d'établir que 
\begin{enumerate} 
\item les composantes irréductibles de $\overline{V}_j \times_S S'$
 sont les schémas $\overline{V}'_{jab}$, et que 
\item la multiplicité de $\overline{V}'_{jab}$ dans 
$\overline{V}_j \times_S S'$ est égale à $n_a m_{jab}$. 
\end{enumerate} 
Remarquons que $V_j \to \overline{V}_j$ est un morphisme birationnel 
de schémas intègres. Les morphismes $V_j \times_S S' \to V_j$
 et $\overline{V}_j \times_S S' \to \overline{V}_j$ sont plats
 et envoient donc les points génériques des composantes irréductibles sur les
 points génériques (uniques) de $V_j$ et de $\overline{V}_j$.
 Il s'ensuit que $V_j \times_S S' \to \overline{V}_j \times_S S'$
 est un morphisme birationnel, donc induit une bijection entre les 
composantes irréductibles et identifie leurs multiplicités.
 Il suffit par conséquent de démontrer que les composantes irréductibles de 
$V_j \times_S S'$ sont les schémas $V'_{jab}$ et 
que la multiplicité de $V'_{jab}$ dans $V_j \times_S S'$
 est égale à $n_a m_{jab}$. Mais cela revient précisément à dire
 que le diagramme 
$$
\xymatrix{
Z_r(V_j) \ar[r] & Z_{r + c}(V_j \times_S S') \\
Z_0(\Spec(\kappa(y))) \ar[r] \ar[u] &
Z_c(\Spec(\kappa(y)) \times_S S') \ar[u]
}
$$
 est commutatif, où les flèches horizontales sont données par changement de base le long de 
$\Spec(\kappa(y)) \times_S S' \to \Spec(\kappa(y))$ et 
les flèches verticales sont les images inverses plates. Cela a été démontré 
dans Homologie de Chow, lemme \ref{chow-lemma-pullback-base-change-pullback}. 

\medskip\noindent
 L'assertion du lemme concernant l'image inverse plate signifie ce qui suit. 
Soient $(S, \delta)$, $X \to Y$, $\alpha$ et $\beta$ comme dans la
 construction de $\alpha \cap \beta$ ci-dessus. Soit $Y' \to Y$ un morphisme
 plat, localement de type fini et de dimension relative $c$. Nous pouvons
 alors prendre pour $\alpha'$ le changement de base de $\alpha$ par $Y' \to Y$
 et pour $\beta'$ l'image inverse plate de $\beta$. La compatibilité avec l'image inverse
 plate signifie que $\alpha' \cap \beta'$ est l'image inverse plate de 
$\alpha \cap \beta$ par $X \times_Y Y' \to Y$. Il s'agit en fait d'un
 cas particulier de la discussion précédente si nous posons $S = Y$ et $S' = Y'$. 
\end{proof} 

\begin{lemma} 
\label{lemma-action-coherent} 
Soient $(S, \delta)$ et $f : X \to Y$ comme ci-dessus. 
Soit $\mathcal{F}$ un $\mathcal{O}_X$-module cohérent
 tel que $\dim(\text{Supp}(\mathcal{F}_y)) \leq r$ pour tout $y \in Y$.
 Soit $\mathcal{G}$ un $\mathcal{O}_Y$-module cohérent
 tel que $\dim_\delta(\text{Supp}(\mathcal{G})) \leq e$. 
Posons $\alpha = [\mathcal{F}/X/Y]_r$
 (exemple \ref{example-family-associated-module}) et 
$\beta = [\mathcal{G}]_e$ (Homologie de Chow, définition 
\ref{chow-definition-cycle-associated-to-coherent-sheaf}). 
Si $\mathcal{F}$ est plat sur $Y$, alors $\alpha \cap \beta =
[\mathcal{F} \otimes_{\mathcal{O}_X} f^*\mathcal{G}]_{r + e}$. 
\end{lemma} 

\begin{proof} 
Observons que 
$$
\text{Supp}(\mathcal{F} \otimes_{\mathcal{O}_X} f^*\mathcal{G}) =
\text{Supp}(\mathcal{F}) \cap f^{-1}\text{Supp}(\mathcal{G}) =
\bigcup\nolimits_{y \in \text{Supp}(\mathcal{G})} \text{Supp}(\mathcal{F}_y)
$$
 Il en résulte qu'il s'agit d'un fermé dont la dimension $\delta$ est $\leq r + e$.
 L'expression 
$[\mathcal{F} \otimes_{\mathcal{O}_X} f^*\mathcal{G}]_{r + e}$
 a donc un sens. 

\medskip\noindent
 Nous utiliserons les notations 
$\beta = \sum n_i[Z_i]$, $y_i \in Z_i$, 
$\alpha_{y_i} = \sum m_{ij} [V_{ij}]$ et 
$\overline{V}_{ij}$ introduites dans la construction 
de $\alpha \cap \beta$. Puisque $\beta = [\mathcal{G}]_e$
 , les $Z_i$ sont les composantes irréductibles de 
$\text{Supp}(\mathcal{G})$ qui sont de dimension $\delta$ égale à $e$.
 De même, les $V_{ij}$ sont les composantes irréductibles 
de $\text{Supp}(\mathcal{F}_{y_i})$ de dimension $r$.
 Il résulte de ceci et de l'égalité du premier 
paragraphe que les $\overline{V}_{ij}$ sont les composantes 
irréductibles de 
$\text{Supp}(\mathcal{F} \otimes_{\mathcal{O}_X} f^*\mathcal{G})$
 de dimension $\delta$ égale à $r + e$.
 Pour démontrer le lemme, il suffit donc désormais de montrer que 
$$
\text{longueur}_{\mathcal{O}_{X, \xi_{ij}}}(
(\mathcal{F} \otimes_{\mathcal{O}_X} f^*\mathcal{G})_{\xi_{ij}})
=
\text{longueur}_{\mathcal{O}_{X_{y_i}, \xi_{ij}}}((\mathcal{F}_{y_i})_{\xi_{ij}})
\cdot
\text{longueur}_{\mathcal{O}_{Y, y_i}}(\mathcal{G}_{y_i})
$$
 D'après le premier paragraphe de la démonstration, le membre de gauche est 
égal à la longueur du $B = \mathcal{O}_{X, \xi_{ij}}$-module 
$$
\mathcal{G}_{y_i}
\otimes_{\mathcal{O}_{Y, y_i}}
\mathcal{F}_{\xi_{ij}} =
M \otimes_A N
$$
 Ici, $M = \mathcal{G}_{y_i}$ est un 
$A = \mathcal{O}_{Y, y_i}$-module de longueur finie, et $N = \mathcal{F}_{\xi_{ij}}$
 est un $B$-module fini tel que $N/\mathfrak m_AN$ soit de longueur finie. 
Puisque $\mathcal{F}$ est plat sur $Y$, le module $N$ est plat sur $A$. Le 
membre de droite de la formule est égal à 
$$
\text{longueur}_B(N/\mathfrak m_A N) \cdot \text{longueur}_A(M)
$$
 Les membres de droite et de gauche de la formule sont donc additifs en $M$
 (on utilise la platitude de $N$ sur $A$). Il suffit ainsi de démontrer la 
formule lorsque $M = \kappa_A$ est le corps résiduel, auquel 
cas elle est immédiate. 
\end{proof} 

\begin{lemma} 
\label{lemma-action-closed} 
Soient $(S, \delta)$ et $f : X \to Y$ comme ci-dessus. Soit $Z \subset X$
 un sous-schéma fermé de dimension relative $\leq r$ sur $Y$. 
Posons $\alpha = [Z/X/Y]_r$ (exemple \ref{example-family-associated-closed}). 
Soit $W \subset Y$ un sous-schéma fermé dont la dimension $\delta$ est $\leq e$. 
Posons $\beta = [W]_e$ (Homologie de Chow, définition 
\ref{chow-definition-cycle-associated-to-closed-subscheme}). 
Si $Z$ est plat sur $Y$, alors $\alpha \cap \beta = [Z \times_Y W]_{r + e}$. 
\end{lemma} 

\begin{proof} 
C'est un cas particulier du lemme \ref{lemma-action-coherent} si
 l'on prend $\mathcal{F} = \mathcal{O}_Z$ et $\mathcal{F} = \mathcal{O}_W$. 
\end{proof} 

\begin{lemma} 
\label{lemma-flat-pullback-as-action} 
Soient $(S, \delta)$ et $f : X \to Y$ comme ci-dessus. Supposons que 
$f$ soit plat de dimension relative $e$. Pour $\beta \in Z_r(Y)$
 , nous avons $f^*\beta = [X/X/Y]_e \cap \beta$ dans $Z_{e + r}(X)$. 
\end{lemma} 

\begin{proof} 
Il suffit de démontrer l'égalité localement sur $X$. Nous pouvons donc supposer 
$Y$ affine. Dans ce cas, $\beta$ est une combinaison linéaire finie à coefficients
 entiers de cycles premiers. Puisque les deux membres de l'égalité sont additifs en $\beta$
 , nous pouvons supposer $\beta = [W]$, où $W$ est un sous-schéma fermé intègre de 
$Y$ de dimension $\delta$ égale à $r$. Soit $f^{-1}(W) = W \times_Y X$ l'image
 inverse schématique de $W$ dans $X$. D'après
 le lemme \ref{lemma-action-closed}, nous avons 
$[X/X/Y]_e \cap \beta = [W \times_Y X]_{r + e}$ et, d'après
 Homologie de Chow, définition \ref{chow-definition-flat-pullback}, 
nous avons $f^*\beta = [f^{-1}(W)]_{r + e}$. Nous obtenons ainsi l'égalité
 voulue. 
\end{proof} 

\begin{lemma} 
\label{lemma-action-push-pull} 
Soit $(S, \delta)$ comme ci-dessus. Soit 
$$
\xymatrix{
X' \ar[r]_f \ar[d] & X \ar[d] \\
Y' \ar[r]^g & Y
}
$$
 un diagramme cartésien de schémas localement de type fini sur $S$
 tel que $g$ soit propre. Soient $r, e \geq 0$. Soit $\alpha$ une famille de 
$r$-cycles sur les fibres de $X/Y$. Soit $\beta' \in Z_e(Y')$.
 Alors $f_*(g^*\alpha \cap \beta') = \alpha \cap g_*\beta'$. 
\end{lemma} 

\begin{proof} 
Puisque nous démontrons une égalité de cycles sur $X$, nous pouvons travailler localement
 sur $Y$; voir le lemme \ref{lemma-action-opens}. Nous pouvons donc supposer $Y$
 affine. Alors $Y'$ est quasi-compact. En particulier, $\beta'$
 est une combinaison linéaire finie de cycles premiers. 
Puisque $- \cap -$ est linéaire en la seconde variable
 (lemme \ref{lemma-action-bilinear}), il suffit de 
démontrer l'égalité lorsque $\beta' = [Z']$ pour un sous-schéma fermé intègre 
$Z' \subset Y'$ de dimension $\delta$ égale à $e$. Posons $Z = g(Z')$. C'est 
un sous-schéma fermé intègre de $Y$ dont la dimension $\delta$ est $\leq e$.
 Pour simplifier, nous supposerons que $Z$ est de dimension $\delta$ égale
 à $e$ et laisserons au lecteur l'autre cas (qui est plus facile). 
Soient $y \in Z$ et $y' \in Z'$ les points génériques. 
Écrivons $\alpha_y = \sum m_j[V_j]$, où $V_j \subset X_y$
 sont des sous-schémas fermés intègres de dimension $r$. 

\medskip\noindent
 Supposons d'abord que $g$ soit une immersion fermée. Alors $g_*\beta' = [Z]$
 et $(g^*\alpha)_{y'} = \sum n_j[V_j]$; cette expression a un sens car 
$V_j$ est contenu dans le sous-schéma fermé $X'_{y'}$ de $X_y$.
 Dans ce cas, l'égalité est donc évidente : dans les deux membres, 
nous obtenons $\sum m_j[\overline{V}_j]$, où $\overline{V}_j$
 est l'adhérence de $V_j$ dans le sous-schéma fermé $X' \subset X$. 

\medskip\noindent
 Revenons au cas général où $\beta' = [Z']$ comme ci-dessus. 
Posons $W = Z \times_Y X$ et $W' = Z' \times_{Y'} X'$.
 Considérons les carrés cartésiens 
$$
\xymatrix{
W \ar[r] \ar[d] & X \ar[d] \\
Z \ar[r] & Y
}
\quad
\xymatrix{
W' \ar[r] \ar[d] & X' \ar[d] \\
Z' \ar[r] & Y'
}
\quad
\xymatrix{
W' \ar[r] \ar[d] & W \ar[d] \\
Z' \ar[r] & Z
}
$$
 Puisque le résultat est connu pour les deux premiers carrés d'après
 le paragraphe précédent, un argument formel montre qu'il suffit 
de démontrer le résultat pour le dernier carré et l'élément 
$\beta' = [Z'] \in Z_e(Z')$. Cela nous ramène au cas traité au 
paragraphe suivant. 

\medskip\noindent
 Supposons que $Y' \to Y$ soit un morphisme génériquement fini entre schémas intègres de dimension
 $\delta$ égale à $e$ et que $\beta' = [Y']$. Dans ce cas, 
$f_*(g^*\alpha \cap \beta')$ et $\alpha \cap g_*\beta'$ sont tous 
deux des cycles qui s'écrivent comme des sommes de cycles premiers dominant $Y$.
 Nous pouvons donc remplacer $Y$ par un sous-schéma ouvert non vide afin de 
vérifier l'égalité. Après un tel remplacement, nous pouvons supposer $g$ fini
 et plat, disons de degré $d \geq 1$. Bien entendu, cela signifie
 que $g_*\beta' = g_*[Y'] = d[Y]$. De plus, $\beta' = [Y'] = g^*[Y]$. Ainsi 
$$
f_*(g^*\alpha \cap \beta') =
f_*(g^*\alpha \cap g^*[Y]) =
f_*f^*(\alpha \cap [Y]) =
d (\alpha \cap [Y]) =
\alpha \cap g_*\beta'
$$
 comme souhaité. La deuxième égalité est le lemme \ref{lemma-action-base-change}
 et la troisième est Homologie de Chow, lemme \ref{chow-lemma-finite-flat}. 
\end{proof}      







\section{Action sur les groupes de Chow} 
\label{section-action-chow} 

\noindent
 Lorsque $\alpha$ est un $r$-cycle relatif, l'opération $\alpha \cap -$
 de la section \ref{section-action} passe au quotient par l'équivalence rationnelle
 et définit une classe bivariante. 

\begin{lemma} 
\label{lemma-closed-in-X-gysin} 
Soit $(S, \delta)$ comme dans la section \ref{section-action}. 
Soit $f : X' \to X$ un morphisme propre de schémas 
localement de type fini sur $S$.
 Soit $(\mathcal{L}, s, i : D \to X)$ comme dans
 Homologie de Chow, définition \ref{chow-definition-gysin-homomorphism}. 
Formons le diagramme 
$$
\xymatrix{
D' \ar[d]_g \ar[r]_{i'} & X' \ar[d]^f \\
D \ar[r]^i & X
}
$$
 comme dans Homologie de Chow, remarque \ref{chow-remark-pullback-pairs}. 
Si $\mathcal{L}|_D \cong \mathcal{O}_D$, alors 
$i^*f_*\alpha' = g_*(i')^*\alpha'$ dans $Z_k(D)$
 pour tout $\alpha' \in Z_{k + 1}(X')$. 
\end{lemma} 

\begin{proof} 
L'énoncé a un sens puisque toutes les opérations sont définies au niveau des cycles ; pour les 
applications de Gysin, voir Homologie de Chow, remarque \ref{chow-remark-gysin-on-cycles}.
 Supposons
 que $\alpha = [W']$ pour un sous-schéma fermé intègre 
$W' \subset X'$. Posons $W = f(W') \subset X$. Si $W' \not \subset D'$,
 alors $W \not \subset D$ et nous voyons que 
$$
[W' \cap D']_k = \text{div}_{\mathcal{L}'|_{W'}}({s'|_{W'}})
\quad\text{et}\quad
[W \cap D]_k = \text{div}_{\mathcal{L}|_W}(s|_W)
$$
 et donc l'image du premier cycle par $f_*$ est égale au second d'après Homologie
 de Chow, lemme \ref{chow-lemma-equal-c1-as-cycles}. L'égalité est
 donc vraie au niveau des cycles. Si $W' \subset D'$, alors 
$W \subset D$ et les deux membres sont nuls par construction. 
\end{proof} 

\begin{lemma} 
\label{lemma-action-gysin} 
Soit $(S, \delta)$ comme dans la section \ref{section-action}. 
Soit $X \to Y$ un morphisme de schémas 
localement de type fini sur $S$. Soit $r \geq 0$ et soit 
$\alpha \in z(X/Y, r)$ un $r$-cycle relatif sur $X/Y$.
 Soit $(\mathcal{L}, s, i : D \to Y)$ comme dans
 Homologie de Chow, définition \ref{chow-definition-gysin-homomorphism}. 
Formons le diagramme cartésien 
$$
\xymatrix{
E \ar[d] \ar[r]_j & X \ar[d] \\
D \ar[r]^i & Y
}
$$
 Voir Homologie de Chow, remarque \ref{chow-remark-pullback-pairs}. 
Si $\mathcal{L}|_D \cong \mathcal{O}_D$, alors, pour $e \in \mathbf{Z}$
 , le diagramme 
$$
\xymatrix{
Z_e(D) \ar[rr]_{i^*\alpha \cap -} & &
Z_{e + r}(E) \\
Z_{e + 1}(Y) \ar[u]^{i^*} \ar[rr]^{\alpha \cap -} & &
Z_{r + e + 1}(X) \ar[u]_{j^*}
}
$$
 est commutatif, les flèches verticales $i^*$ et $j^*$ étant les 
applications de Gysin sur les cycles définies
 dans Homologie de Chow, remarque \ref{chow-remark-gysin-on-cycles}. 
\end{lemma} 

\begin{proof} 
Remarque préliminaire.
 Supposons que $g : Y' \to Y$ soit une enveloppe (Homologie de Chow, définition 
\ref{chow-definition-envelope}). 
Notons $D', i', E', j', X', \alpha'$ les changements de base de 
$D, i, E, j, X, \alpha$ par $g$, et notons $f : X' \to X$ la projection. 
Supposons le lemme vrai pour $D', i', E', j', X', Y', \alpha'$.
 Alors, si $\beta' \in Z_{e + 1}(Y')$, nous avons 
\begin{align*} 
i^*\alpha \cap i^*g_*\beta'
 & =
 i^*\alpha \cap f_*(i')^*\beta' \\
 & =
 f_*(f^*i^*\alpha \cap (i')^*\beta') \\ 
& =
 f_*((i')^*\alpha' \cap (i')^*\beta') \\ 
& =
 f_*((j')^*(\alpha' \cap \beta')) \\ 
& =
 j^*(f_*(f^*\alpha \cap \beta')) \\ 
& =
 j^*(\alpha \cap g_*\beta') 
\end{align*} 
Ici, la 
première égalité est le lemme \ref{lemma-closed-in-X-gysin},
 la deuxième le lemme \ref{lemma-action-push-pull},
 la troisième la définition de $\alpha'$,
 la quatrième résulte de l'hypothèse selon laquelle notre lemme est vrai pour 
$D', i', E', j', X', \alpha'$,
 la cinquième est le lemme \ref{lemma-closed-in-X-gysin},
 et la sixième le lemme \ref{lemma-action-push-pull}. Nous voyons
 ainsi que notre lemme est vrai pour tout cycle appartenant à l'image de 
$g_* : Z_{e + 1}(Y') \to Z_e(Y)$. Cependant, puisque $g$ est complètement 
décomposé, cette application est surjective ; nous en concluons que le lemme est vrai 
pour $D, i, E, j, X, Y, \alpha$. 

\medskip\noindent
 Soit $\beta \in Z_{e + 1}(Y)$. Nous devons montrer que 
$(D \to Y)^*\alpha \cap i^*\beta = j^*(\alpha \cap \beta)$
 au niveau des cycles sur $E$. Cette question est locale sur $E$ ; nous pouvons
 donc remplacer $X$ et $Y$ par des sous-schémas ouverts. (Nous utilisons ici le fait que la 
formation des opérateurs $i^*$, $j^*$, $\alpha \cap - $ et 
$(D \to Y)^*\alpha \cap -$ commute à la localisation. Cela est 
évident pour les applications de Gysin et résulte du
 lemme \ref{lemma-action-opens} pour les autres.) Nous
 pouvons donc supposer que $X$ et $Y$ sont affines
 et nous ramener au cas étudié dans le paragraphe suivant. 

\medskip\noindent
 Supposons $X$ et $Y$ quasi-compacts. D'après le premier paragraphe de la preuve 
et le lemme \ref{lemma-get-cycles}, nous pouvons en outre supposer que $\alpha$
 appartient à l'image de (\ref{equation-cycle-classes}). Par linéarité 
des opérations considérées, nous pouvons supposer que $\alpha = [Z/X/Y]_r$
 pour un sous-schéma fermé $Z \subset X$, plat et de dimension 
relative $\leq r$ sur $Y$. De plus, puisque $Y$ est quasi-compact, le cycle 
$\beta$ est une combinaison linéaire finie de cycles premiers. Les opérations 
considérées étant linéaires, il suffit de
 démontrer l'égalité lorsque $\beta = [W]$ pour un sous-schéma fermé intègre 
$W \subset Y$ de $\delta$-dimension $e + 1$. 

\medskip\noindent
 Si $W \subset D$, alors, d'une part, $i^*[W] = 0$ et, d'autre
 part,$\alpha \cap [W]$ est supporté sur $E$, de sorte que 
$j^*(\alpha \cap [W]) = 0$. L'égalité est donc vraie dans ce cas. 

\medskip\noindent
 Supposons $W \not \subset D$. Alors $i^*[W] = [D \cap W]_e$.
 Notons que l'image inverse $i^*\alpha$ de $\alpha = [Z/X/Y]_r$
 par $i$ est $[(E \cap Z)/E/D]_r$ et que 
$(E \cap Z) = E \times_Y Z = D \times_Y Z$
 est plat sur $D$. Le lemme \ref{lemma-action-closed}, appliqué
 deux fois, donne donc 
$$
i^*\alpha \cap i^*[W] =
[(E \cap Z) \times_D (D \cap W)]_{r + e} =
[E \cap (Z \times_Y W)]_{r + e} =
j^*(\alpha \cap [W])
$$
 comme voulu. 
\end{proof} 

\begin{proposition} 
\label{proposition-get-bivariant-class} 
Soit $(S, \delta)$ comme dans la section \ref{section-action}. Soit $X \to Y$
 un morphisme de schémas localement de type fini sur $S$. Soit $r \geq 0$
 et soit $\alpha \in z(X/Y, r)$ un $r$-cycle relatif sur $X/Y$.
 La règle qui, à tout morphisme $g : Y' \to Y$ localement de type fini
 et à tout $e \in \mathbf{Z}$, associe l'opération 
$$
g^*\alpha \cap - : Z_e(Y') \to Z_{r + e}(X')
$$
 où $X' = Y' \times_Y X$, passe au quotient par l'équivalence rationnelle 
et définit une classe bivariante $c(\alpha) \in A^{-r}(X \to Y)$. 
\end{proposition} 

\begin{proof} 
L'opération passe au quotient par l'équivalence rationnelle d'après 
le lemme \ref{lemma-action-gysin} et 
Homologie de Chow, lemme \ref{chow-lemma-factors-through-rational-equivalence}. L'opération
 ainsi obtenue sur les groupes de Chow est une classe bivariante d'après 
Homologie de Chow, lemme \ref{chow-lemma-bivariant-weaker} et 
les
 lemmes \ref{lemma-action-push-pull}, \ref{lemma-action-base-change}, et 
\ref{lemma-action-gysin}. 
\end{proof} 

\begin{remark} 
\label{remark-characterize-relative-cycles} 
Soit $(S, \delta)$ comme dans la section \ref{section-action}. Soit $X \to Y$
 un morphisme de schémas localement de type fini sur $S$. Soit $r \geq 0$.
 Soit $c$ une règle qui, à tout morphisme $g : Y' \to Y$ localement de type
 fini et à tout $e \in \mathbf{Z}$, associe une opération 
$$
c \cap - : Z_e(Y') \to Z_{r + e}(X')
$$
 compatible avec l'image directe propre et l'image inverse plate, ainsi qu'avec les applications de Gysin comme dans
 le lemme \ref{lemma-action-gysin}. Nous affirmons qu'il existe un $r$-cycle relatif 
$\alpha$ sur $X/Y$ tel que $c \cap = g^*\alpha \cap -$ pour tout $g$ comme ci-dessus.
 Si nous en avons besoin, nous en donnerons ici un énoncé précis et une preuve détaillée. 
\end{remark} 









\section{Composition des familles de cycles sur les fibres} 
\label{section-compose-families} 

\noindent
 Soient $X \to Y \to S$ des morphismes de schémas, tous deux localement de type fini. 
Soient $r, e \geq 0$. Soit $\alpha$ une famille de $r$-cycles 
sur les fibres de  $X/Y$  et  $\beta$  une famille de  $e$-cycles sur les
 fibres de  $Y/S$. Alors nous obtenons une famille 
de  $(r + e)$-cycles $\alpha \circ \beta$  sur les fibres de  $X/S$
  en posant 
$$
(\alpha \circ \beta)_s = (Y_s \to Y)^*\alpha \cap \beta_s
$$
 Plus précisément, l'expression $(Y_s \to Y)^*\alpha$ désigne 
le changement de base de $\alpha$  par  $Y_s \to Y$  vers une famille de  $r$-cycles 
sur les fibres de  $X_s/Y_s$  et l'opération  $- \cap -$
  a été définie et étudiée dans la section \ref{section-action}\footnote{Précisément, 
nous prenons $s = \Spec(\kappa(s))$ comme schéma de base, avec $\delta(s) = 0$.} 

\begin{lemma} 
\label{lemma-construction-bilinear} La 
construction ci-dessus est bilinéaire, c'est-à-dire que nous avons  
$(\alpha_1 + \alpha_2) \circ \beta \alpha_1 \circ \beta +
\alpha_1 \circ \beta$ et  $\alpha \circ (\beta_1 + \beta_2) =
\alpha \circ \beta_1 + \alpha \circ \beta_2$. 
\end{lemma} 

\begin{proof} 
Démonstration omise. Indication : sur les fibres, la construction est bilinéaire 
d'après le lemme \ref{lemma-action-bilinear}. 
\end{proof} 

\begin{lemma} 
\label{lemma-construction-opens} 
Si  $U \subset X$  et  $V \subset Y$  sont ouverts et  $f(U) \subset V$, alors  
$(\alpha \circ \beta)|_U$ est égal à  $\alpha|_U \circ \beta|_V$. 
\end{lemma} 

\begin{proof} 
Démonstration omise. Indication : sur les fibres, appliquer
 le lemme \ref{lemma-action-opens}. 
\end{proof} 

\begin{lemma} 
\label{lemma-construction-base-change} 
La formation de  $\alpha \circ \beta$  est compatible avec le changement de base. 
\end{lemma} 

\begin{proof} 
Soit  $g : S' \to S$  un morphisme de schémas. 
Notons  $X' \to Y'$  le changement de base de  $X \to Y$  par  $g$. Notons  
$\alpha'$  le changement de base de  $\alpha$  par rapport à  $Y' \to Y$.
 Notons  $\beta'$  le changement de base de  $\beta$  par rapport à  $S' \to S$.
 L'assertion signifie que $\alpha' \circ \beta'$  est le changement de base 
de  $\alpha \circ \beta$  par  $g : S' \to S$. 

\medskip\noindent
 Soit  $s' \in S'$  un point d'image  $s \in S$. Alors 
$$
(\alpha' \circ \beta')_{s'} = (Y'_{s'} \to Y')^*\alpha' \cap \beta'_{s'}
$$
 Nous observons que 
$$
(Y'_{s'} \to Y')^*\alpha' =
(Y'_{s'} \to Y')^*(Y' \to Y)^*\alpha =
(Y'_{s'} \to Y_s)^*(Y_s \to Y)^*\alpha
$$
 et que  $\beta'_{s'}$  est le changement de base de  $\beta_s$  par  
$s' = \Spec(\kappa(s')) \to \Spec(\kappa(s)) = s$.
 Par conséquent, le résultat découle du lemme \ref{lemma-action-base-change} 
appliqué à $(Y_s \to Y)^*\alpha$, $\beta_s$,  
$X_s \to Y_s \to s$, et au changement de base par  $s' \to s$. 
\end{proof} 

\begin{lemma} 
\label{lemma-construction-coherent} 
Soient  $f : X \to Y$  et  $Y \to S$  des morphismes de schémas, tous deux localement 
de type fini. Soient $r, e \geq 0$. Soit $\mathcal{F}$ un 
$\mathcal{O}_X$-module quasi-cohérent de type fini, avec  
$\dim(\text{Supp}(\mathcal{F}_y)) \leq r$  pour tout $y \in Y$. 
Soit  $\mathcal{G}$  un  $\mathcal{O}_Y$-module quasi-cohérent de 
type fini, avec  $\dim(\text{Supp}(\mathcal{G}_s)) \leq e$  pour tout  $s \in S$. 
Si  $\alpha = [\mathcal{F}/X/Y]_r$  et  $\beta = [\mathcal{G}/Y/S]_e$
 (Exemple  \ref{example-family-associated-module}) et  $\mathcal{F}$
  est plat sur  $Y$, alors $\alpha \circ \beta =
[\mathcal{F} \otimes_{\mathcal{O}_X} f^*\mathcal{G}/X/S]_{r + e}$. 
\end{lemma} 

\begin{proof} 
Tout d'abord, nous observons que  $\mathcal{F} \otimes_{\mathcal{O}_X} f^*\mathcal{G}$
  est un  $\mathcal{O}_X$-module quasi-cohérent de type fini. 
Soit  $s \in S$. Observons que 
$$
(\mathcal{F} \otimes_{\mathcal{O}_X} f^*\mathcal{G})_s =
\mathcal{F}_s \otimes_{\mathcal{O}_{X_s}} f_s^*\mathcal{G}_s
$$
 par l'exactitude à droite du produit tensoriel. De plus, $\mathcal{F}_s$
 est plat sur $Y_s$, puisqu'il s'obtient par changement de base d'un module plat. Ainsi, l'égalité 
$(\alpha \circ \beta)_s =
[(\mathcal{F} \otimes_{\mathcal{O}_X} f^*\mathcal{G})_s]_{r + e}$
 découle du lemme \ref{lemma-action-coherent}. 
\end{proof} 

\begin{lemma} 
\label{lemma-construction-closed} 
Soient  $f : X \to Y$  et  $Y \to S$  des morphismes de schémas, tous deux localement 
de type fini. Soient $r, e \geq 0$. Soit $Z \subset X$ un sous-schéma
 fermé de dimension relative  $\leq r$  sur  $Y$. 
Soit  $W \subset Y$  un sous-schéma fermé de dimension relative  $\leq e$
  sur  $S$. Si $\alpha = [Z/X/Y]_r$  et  $\beta = [W/Y/S]_e$
  (Exemple  \ref{example-family-associated-closed}) et  $Z$  est plat sur $Y$, 
alors  $\alpha \circ \beta = [Z \times_Y W/X/S]_{r + e}$. 
\end{lemma} 

\begin{proof} 
Il s'agit d'un cas particulier du lemme \ref{lemma-construction-coherent} si 
nous prenons $\mathcal{F} = \mathcal{O}_Z$ et  $\mathcal{F} = \mathcal{O}_W$. 
\end{proof} 

\begin{lemma} 
\label{lemma-flat-pullback-as-composition} Soient
  $f : X \to Y$  et  $Y \to S$  des morphismes de schémas, tous deux localement de 
type fini. Supposons que  $f$  soit plat de dimension relative  $e$. Soit  $\beta$  une
 famille de  $r$-cycles sur les fibres de  $Y/S$. Alors  
$f^*\beta = [X/X/Y]_e \circ \beta$
  comme familles de $(e + r)$-cycles sur  $X/S$. 
\end{lemma} 

\begin{proof} 
En déroulant les définitions et en utilisant la compatibilité de la formation de 
$[X/X/Y]_e$ avec le changement de base 
(lemme \ref{lemma-family-associated-closed}), on conclut par le 
lemme \ref{lemma-flat-pullback-as-action}. 
\end{proof} 

\begin{lemma} 
\label{lemma-construction-push-pull} 
Soit  $S$  un schéma. Soit 
$$
\xymatrix{
X' \ar[r]_f \ar[d] & X \ar[d] \\
Y' \ar[r]^g & Y
}
$$
 un diagramme cartésien de schémas localement de type fini sur  $S$
  avec $g$ propre. Soient $r, e \geq 0$. Soit $\alpha$ une famille de 
$r$-cycles sur les fibres de  $X/Y$. Soit  $\beta'$  une famille de  
$e$-cycles sur les fibres de  $Y'/S$. Alors nous avons 
$f_*(g^*(\alpha) \circ \beta') = \alpha \circ g_*\beta'$. 
\end{lemma} 

\begin{proof}
 En déroulant les définitions, cela découle du 
lemme \ref{lemma-action-push-pull}. 
\end{proof} 

\begin{lemma} 
\label{lemma-construction-composition} 
Soit $(S, \delta)$ comme dans Homologie de Chow, situation \ref{chow-situation-setup}. 
Soient $X \to Y \to Z$  des morphismes de schémas localement de type fini sur  $S$. 
Soient $r, s, e \geq 0$. Alors 
$$
(\alpha \circ \beta) \cap \gamma = \alpha \cap (\beta \cap \gamma)
\quad\text{dans}\quad Z_{r + s + e}(X)
$$
 où  $\alpha$  est une famille de  $r$-cycles sur les fibres de $X/Y$,  
$\beta$  est une famille de  $s$-cycles sur les fibres de  $Y/Z$, et  $\gamma \in Z_e(Z)$. 
\end{lemma} 

\begin{proof} 
Puisque nous prouvons une égalité de cycles sur  $X$, nous pouvons travailler localement
 sur $Z$ ; voir le lemme \ref{lemma-action-opens}. Ainsi nous pouvons supposer que  $Z$
  est affine. En particulier,  $\gamma$  est une combinaison linéaire finie de cycles 
premiers. Puisque $- \cap -$ est linéaire en la seconde variable
 (lemme \ref{lemma-action-bilinear}), il suffit de
 prouver l'égalité lorsque $\gamma = [W]$ pour un sous-schéma fermé intègre 
$W \subset Z$ de $\delta$-dimension $e$. 

\medskip\noindent
 Soit $z \in W$ le point générique. Écrivons $\beta_z = \sum m_j[V_j]$
 dans $Z_s(Y_z)$. Alors  $\beta \cap \gamma$  est égal à  $\sum m_j[\overline{V}_j]$
  où  $\overline{V}_j \subset Y$  est un sous-schéma fermé intègre dont 
l'image par $Y \to Z$ est contenue dans $W$ et dont la fibre générique est $V_j$.
 Soit $y_j \in V_j$ le point générique. Nous le considérons aussi
 comme le point générique de $\overline{V}_j$ (envoyé sur $z$ dans $W$). 
Écrivons $\alpha_{y_j} = \sum n_{jk} [W_{jk}]$ dans $Z_r(X_{y_j})$.
 Alors $\alpha \cap (\beta \cap \gamma)$  est égal à 
$$
\sum m_j n_{jk} [\overline{W}_{jk}]
$$
 où $\overline{W}_{jk} \subset X$ est un sous-schéma fermé intègre dont 
l'image par $X \to Y$ est contenue dans $\overline{V}_j$ et dont la fibre générique est $W_{jk}$. 

\medskip\noindent
 D'autre part, considérons 
$$
(\alpha \circ \beta)_z = (Y_z \to Y)^*\alpha \cap \beta_z =
(Y_z \to Y)^*\alpha \cap (\sum m_j [V_j])
$$
 Par la construction de $- \cap -$, ceci est égal au cycle 
$$
\sum m_j n_{jk} [(\overline{W}_{jk})_z]
$$
 sur  $X_z$. Ainsi, par définition, nous obtenons 
$$
(\alpha \circ \beta) \cap [W] =
\sum m_j n_{jk} [\widetilde{W}_{jk}]
$$
 où $\widetilde{W}_{jk} \subset X$ est un sous-schéma fermé intègre dont
 l'image par $X \to Z$ est contenue dans $W$ et dont la fibre générique est $(\overline{W}_{jk})_z$.
 Il est clair que $\widetilde{W}_{jk} = \overline{W}_{jk}$
 , ce qui achève la preuve. 
\end{proof} 









\section{Composition des cycles relatifs} 
\label{section-compose} 

\noindent
 Soit  $S$  un schéma localement noethérien. Soit  $X \to Y$
  un morphisme de schémas localement de type fini sur  $S$.
 Nous allons définir une application 
$$
z(X/Y, r) \otimes_\mathbf{Z} z(Y/S, e) \longrightarrow z(X/S, r + e),\quad
\alpha \otimes \beta \longmapsto \alpha \circ \beta
$$
 en utilisant la construction de la section \ref{section-compose-families}. 
Nous savons déjà que la construction est bilinéaire
 (lemme \ref{lemma-construction-bilinear}) ; 
pour obtenir la flèche affichée, il reste donc à établir le lemme suivant. 

\begin{lemma} 
\label{lemma-well-defined} 
Si  $\alpha$  et  $\beta$  sont des cycles relatifs, alors il en est de même pour  $\alpha \circ \beta$. 
\end{lemma} 

\begin{proof}
 La formation de $\alpha \circ \beta$ est compatible avec le changement de base d'après
 le lemme \ref{lemma-construction-base-change}. Ainsi, nous pouvons supposer que 
$S$ est le spectre d'un anneau de valuation discrète avec point générique  
$\eta$  et point fermé $0$  et nous devons montrer que  
$sp_{X/S}((\alpha \circ \beta)_\eta) = (\alpha \circ \beta)_0$.
 Puisque nous essayons de prouver une égalité de cycles, nous pouvons travailler
 localement sur  $Y$ et $X$ (on utilise les
 lemmes \ref{lemma-construction-opens} et 
\ref{lemma-specialization-flat-pullback} pour 
voir que les constructions commutent avec la restriction). Ainsi, nous
 pouvons supposer que $X$ et $Y$ sont affines. 
D'après le lemme \ref{lemma-get-cycles}, 
nous pouvons trouver un morphisme propre complètement décomposé  
$g : Y' \to Y$  tel que $g^*\alpha$  soit dans l'image de 
(\ref{equation-cycle-classes}). 

\medskip\noindent
 Comme le morphisme $g_\eta : Y'_\eta \to Y_\eta$
 est complètement décomposé, nous pouvons trouver $\beta'_\eta \in Z_e(Y'_\eta)$ tel que 
$\beta_\eta = \sum g_{\eta, *}\beta'_\eta$ ; voir
 Homologie de Chow, lemme \ref{chow-lemma-envelope}. 
Posons $\beta'_0 = sp_{Y'/S}(\beta'_\eta)$, de sorte que 
$\beta' = (\beta'_\eta, \beta'_0)$ soit un $e$-cycle relatif sur  $Y'/S$. Alors  
$g_*\beta'$  et  $\beta$ sont des $e$-cycles relatifs
 sur $Y/S$ (lemme \ref{lemma-relative-cycle-functoriality})
 qui ont la même valeur en $\eta$ et sont donc égaux 
(lemme \ref{lemma-uniqueness-extension}). Par linéarité 
(lemme \ref{lemma-construction-bilinear}),
 il suffit de montrer que $\alpha \circ g_*\beta'$
 est un $(r + e)$-cycle relatif. 

\medskip\noindent
 Posons $X' = X \times_Y Y'$ et notons $f : X' \to X$ la projection. D'après
 le lemme \ref{lemma-construction-push-pull}, on a 
$\alpha \circ g_*\beta' = f_*(g^*\alpha \circ \beta')$. 
D'après le lemme \ref{lemma-relative-cycle-functoriality}, 
il suffit de montrer que  $g^*\alpha \circ \beta'$
  est un  $(r + e)$-cycle relatif. En utilisant le lemme \ref{lemma-get-cycles-dvr} 
et la bilinéarité, cela nous ramène au cas discuté dans le paragraphe suivant. 

\medskip\noindent
 Supposons  $\alpha = [Z/X/Y]_r$  et  $\beta = [W/Y/S]$
  où  $Z \subset X$  est un sous-schéma fermé plat et de dimension
 relative  $\leq r$  sur  $Y$  et  $W \subset Y$  est un 
sous-schéma fermé plat et de dimension relative  $\leq e$  sur  $S$. 
D'après le lemme \ref{lemma-construction-closed}, on a 
$$
\alpha \circ \beta = [Z \times_X W/X/S]_{r + e}
$$
 et  $Z \times_X W \subset X$  est un sous-schéma fermé plat sur  $S$
  de dimension relative $\leq r + e$. C'est un $(r + e)$-cycle relatif d'après
 le lemme \ref{lemma-family-associated-closed-specialization}. 
\end{proof} 

\begin{lemma} 
\label{lemma-composition-fundamental-cycles} Soient
  $f : X \to Y$  et  $g : Y \to S$  des morphismes de schémas. Supposons
 que $S$ soit localement noethérien, que $g$ soit localement de type fini 
et plat de dimension relative  $e \ge 0$, et  $f$  localement de type fini
 et plat de dimension relative  $r \geq 0$. Alors  
$[X/X/Y]_r \circ [Y/Y/S]_e = [X/X/S]_{r + e}$  dans  $z(X/S, r + e)$. 
\end{lemma} 

\begin{proof} 
C'est un cas particulier du lemme \ref{lemma-construction-closed}. 
\end{proof} 






\section{Comparaison avec Suslin et Voevodsky} 
\label{section-compare} 

\noindent
 Nous avons cherché à employer les mêmes notations que dans \cite{SV}, sauf que notre
 notation pour les cycles est tirée 
d'Homologie de Chow, section \ref{chow-section-cycles} et suivantes.
 Voici la comparaison : 
\begin{enumerate} 
\item Dans  \cite[Section 3.1]{SV} , il existe une notion
 de « cycle relatif », de « cycle relatif de dimension  $r$ » et 
de « cycle relatif équidimensionnel de dimension $r$ ». 
Il n'existe pas de notion correspondante dans ce chapitre. Par conséquent, les groupes  
$Cycl(X/S, r)$,  $Cycl_{equi}(X/S, r)$,  
$PropCycl(X/S, r)$ et  $PropCycl_{equi}(X/S, r)$ 
n'ont pas de correspondants dans ce chapitre. 
\item En bas de  \cite[page 36]{SV} , les groupes  
$z(X/S, r)$,  $c(X/S, r)$,  $z_{equi}(X/S, r)$ et $c_{equi}(X/S, r)$
  sont définis. Ceux-ci correspondent à nos notions lorsque $S$ est séparé et
 noethérien, et $X \to S$ est séparé et de type fini. 
\item Dans  \cite{SV} , le symbole  $z(X/S, r)$  est parfois utilisé pour le 
préfaisceau $S' \mapsto z(S' \times_S X/S', r)$  sur la catégorie des schémas
 de type fini sur  $S$. De même pour  
$c(X/S, r)$,  $z_{equi}(X/S, r)$ et  $c_{equi}(X/S, r)$. 
\item Le changement de base, l'image inverse plate et l'image directe 
propre définis dans \cite{SV} coïncident avec les nôtres lorsqu'ils sont tous deux définis. 
\item Pour $\alpha \in z(X/S, r)$, l'opération 
$\alpha \cap - : Z_e(S) \to Z_{e + r}(X)$ définie dans 
la section \ref{section-action} coïncide avec
 l'opération $Cor(\alpha, -)$ de \cite[Section 3.7]{SV} lorsque
 les deux sont définies. 
\item Pour $X \to Y \to S$, la loi de composition 
$z(X/Y, r) \otimes_\mathbf{Z} z(Y/S, e) \longrightarrow z(X/S, r + e)$
 définie dans la section \ref{section-compose}
 coïncide avec l'opération $Cor_{X/Y}(-, -)$ de 
\cite[Corollary 3.7.5]{SV}. 
\end{enumerate} 











\section{Cycles relatifs dans le cas non noethérien} 
\label{section-non-noetherian} 

\noindent
 Nous invitons le lecteur à sauter cette section. 

\medskip\noindent
 Soit  $f : X \to S$  un morphisme de schémas de présentation finie. 
Soit $r \geq 0$. On note  $Hilb(X/S, r)$  l'ensemble des sous-schémas fermés  
$Z \subset X$ tels que  $Z \to S$  soit plat, de présentation finie,
 et de dimension relative $\leq r$. Nous considérons l'homomorphisme de groupes 
\begin{equation} 
\label{equation-cycle-classes-general} 
\begin{matrix} 
\text{groupe abélien libre} \\ 
\text{sur }Hilb(X/S, r) 
\end{matrix} 
\longrightarrow
\begin{matrix} 
\text{familles de }r\text{-cycles}\\ 
\text{sur les fibres de }X/S 
\end{matrix} 
\end{equation} 
envoyant  $\sum n_i[Z_i]$  vers  $\sum n_i[Z_i/X/S]_r$. 

\begin{lemma} 
\label{lemma-relative-r-cycle-general} 
Soit  $S$  un schéma quasi-compact et quasi-séparé. 
Soit  $f : X \to S$  un morphisme de présentation finie. 
Soit  $r \geq 0$  et soit  $\alpha$  une famille de $r$-cycles sur les fibres de $X/S$. 
Les assertions suivantes sont équivalentes : 
\begin{enumerate} 
\item il existe un diagramme cartésien 
$$
\xymatrix{
X \ar[r] \ar[d] & X_0 \ar[d] \\
S \ar[r] & S_0
}
$$
 où  $X_0 \to S_0$  est un morphisme de type fini de schémas noethériens
 et $\alpha_0 \in z(X_0/S_0, r)$  tel que  $\alpha$  soit le changement de base 
de  $\alpha_0$  par $S \to S_0$
\item il existe un morphisme propre complètement décomposé  $g : S' \to S$
  de présentation finie tel que  $g^*\alpha$  soit dans l'image de 
(\ref{equation-cycle-classes-general}). 
\end{enumerate} 
\end{lemma} 

\begin{proof} 
Supposons donnés un diagramme et $\alpha_0 \in z(X_0/S_0, r)$ comme dans (1). D'après
 le lemme \ref{lemma-get-cycles}, il existe un morphisme propre complètement 
décomposé  $g_0 : S'_0 \to S_0$  tel que $g_0^*\alpha_0$  soit 
dans l'image de (\ref{equation-cycle-classes-general}). En effet, puisque  
$S'_0$ est noethérien, tout sous-schéma fermé de  $S'_0 \times_{S_0} X_0$
  est de présentation finie sur  $S'_0$. En posant  $S' = S \times_{S_0} S'_0$
  et en utilisant le changement de base par $S' \to S'_0$, on voit que (2) est vérifiée. 

\medskip\noindent
 Réciproquement, supposons (2) vérifiée. Choisissons un morphisme propre complètement 
décomposé  $g : S' \to S$  de présentation finie tel que  $g^*\alpha$ soit dans 
l'image de (\ref{equation-cycle-classes-general}). Posons $X' = S' \times_S X$. 
Écrivons $g^*\alpha = \sum n_a [Z_a/X'/S']_r$  pour certains sous-schémas fermés  $Z_a \subset X'$
 , plats, de présentation finie et de dimension relative  
$\leq r$  sur  $S'$. 

\medskip\noindent
 Écrivons $S = \lim S_i$, où le système d'indices est filtrant, les morphismes de transition
 sont affines et les $S_i$ sont de type fini sur $\mathbf{Z}$ 
; voir Limites, proposition \ref{limits-proposition-approximate}.
 Pour $i$ suffisamment grand, il existe 
\begin{enumerate} 
\item un morphisme propre complètement décomposé  $g_i : S'_i \to S_i$
  dont le changement de base à  $S$  est  $g : S' \to S$, 
\item en posant $X'_i = S'_i \times_{S_i} X_i$
 , des sous-schémas fermés $Z_{ai} \subset X'_i$, plats et de
 dimension relative $\leq r$ sur $S'_i$, dont le changement de base à $S'$
 est $Z_a$. 
\end{enumerate} 
Pour ce faire, on utilise Limites, lemmes 
\ref{limits-lemma-descend-finite-presentation}, 
\ref{limits-lemma-descend-closed-immersion-finite-presentation}, 
\ref{limits-lemma-descend-flat-finite-presentation}, 
\ref{limits-lemma-descend-surjective}, 
\ref{limits-lemma-eventually-proper} et 
\ref{limits-lemma-limit-dimension}, ainsi 
que 
Compléments sur les morphismes, lemme 
\ref{more-morphisms-lemma-descend-cd}. 
Considérons $\alpha'_i = \sum n_a [Z_{ai}/X'_i/S_i]_r \in z(X'_i/S'_i, r)$.
 Le changement de base de  $\alpha'_i$  vers une famille de  $r$-cycles sur les fibres
 de $X'/S'$ coïncide avec le changement de base $g^*\alpha$ par construction. 

\medskip\noindent
 Posons $S''_i = S'_i \times_{S_i} S'_i$ et $X''_i = S''_i \times_{S_i} X_i$
 , puis $S'' = S' \times_S S'$ et $X'' = S'' \times_S X$.
 Notons $\text{pr}_1, \text{pr}_2 : S'' \to S'$ et 
$\text{pr}_1, \text{pr}_2 : S''_i \to S'_i$ les projections.
 Les $r$-cycles relatifs  $\text{pr}_1^*\alpha'_i$  et  $\text{pr}_1^*\alpha'_i$
  sur la base  $X''_i/S''_i$ deviennent après changement de base la même famille de $r$-cycles 
sur les fibres de $X''/S''$  parce que  
$\text{pr}_1^*g^*\alpha = \text{pr}_1^*g^*\alpha$. 
Par conséquent, l'image du morphisme $S'' \to S''_i$ est contenue dans $E =
\{s \in S''_i : (\text{pr}_1^*\alpha'_i)_s = (\text{pr}_1^*\alpha'_i)_s\}$. 
D'après le lemme \ref{lemma-relative-cycles-equal}, il s'agit d'un sous-ensemble fermé.
 Comme $S'' = \lim_{i' \geq i} S''_{i'}$, le 
lemme 
\ref{limits-lemma-limit-contained-in-constructible} de Limites montre 
que, pour un certain $i' \geq i$, l'image du morphisme $S''_{i'} \to S''_i$
 est contenue dans $E$. Par conséquent, après avoir remplacé  $i$  par  $i'$, nous pouvons supposer que 
$\text{pr}_1^*\alpha'_i = \text{pr}_1^*\alpha'_i$. 
D'après le lemme \ref{lemma-descend-family}, nous
 obtenons une famille unique  $\alpha_i$
  de  $r$-cycles sur les fibres de  $X_i/S_i$
 avec  $g_i^*\alpha_i = \alpha'_i$  (cela utilise le fait que  $S'_i \to S_i$
  est complètement décomposé). D'après
 le lemme \ref{lemma-relative-cycles-h-descent},
 on a $\alpha_i \in z(X_i/S_i, r)$.
 L'unicité dans le lemme \ref{lemma-descend-family} implique que le 
changement de base de $\alpha_i$ est $\alpha$ ; ainsi (1) est vérifiée. 
\end{proof} 

\noindent
 {\bf Discussion.} 
Si $f : X \to S$, $r$ et $\alpha$ sont comme dans
 le lemme \ref{lemma-relative-r-cycle-general}, il est naturel de 
dire que $\alpha$ est un {\it $r$-cycle relatif sur $X/S$} si
 les conditions équivalentes
 (1) et (2) du lemme \ref{lemma-relative-r-cycle-general} sont satisfaites.
 Cette définition possède de nombreuses bonnes propriétés ; par 
exemple, elle est compatible avec la définition précédente lorsque $S$
 est noethérien, et la plupart des résultats de 
la section \ref{section-families-specialization} se généralisent
 à ce cadre. 

\medskip\noindent
 Nous pouvons encore généraliser comme suit. 
Supposons $S$ arbitraire et $f : X \to S$ localement de présentation finie. 
Soient $r \geq 0$ et $\alpha$ une famille de $r$-cycles, notée $\alpha$, sur les fibres de 
$X/S$. Alors $\alpha$ est un {\it $r$-cycle relatif sur $X/S$} si, pour tous 
ouverts affines $U \subset X$ et $V \subset S$ tels que $f(U) \subset V$
 , la restriction $\alpha|_U$ est un $r$-cycle relatif sur $U/V$ au
 sens du paragraphe précédent. Là encore, beaucoup des résultats antérieurs
 se généralisent à ce cadre. 

\medskip\noindent
 Si nous avons besoin de ces généralisations, nous les énoncerons et les prouverons 
soigneusement ici. 



















\begin{multicols}{2}[\section{Autres chapitres}]
\noindent
Préliminaires
\begin{enumerate}
\item \hyperref[introduction-section-phantom]{Introduction}
\item \hyperref[conventions-section-phantom]{Conventions}
\item \hyperref[sets-section-phantom]{Théorie des ensembles}
\item \hyperref[categories-section-phantom]{Catégories}
\item \hyperref[topology-section-phantom]{Topologie}
\item \hyperref[sheaves-section-phantom]{Faisceaux sur les espaces}
\item \hyperref[sites-section-phantom]{Sites et faisceaux}
\item \hyperref[stacks-section-phantom]{Champs}
\item \hyperref[fields-section-phantom]{Corps}
\item \hyperref[algebra-section-phantom]{Algèbre commutative}
\item \hyperref[brauer-section-phantom]{Groupes de Brauer}
\item \hyperref[homology-section-phantom]{Algèbre homologique}
\item \hyperref[derived-section-phantom]{Catégories dérivées}
\item \hyperref[simplicial-section-phantom]{Méthodes simpliciales}
\item \hyperref[more-algebra-section-phantom]{Compléments d'algèbre}
\item \hyperref[smoothing-section-phantom]{Lissage des morphismes d'anneaux}
\item \hyperref[modules-section-phantom]{Faisceaux de modules}
\item \hyperref[sites-modules-section-phantom]{Modules sur les sites}
\item \hyperref[injectives-section-phantom]{Objets injectifs}
\item \hyperref[cohomology-section-phantom]{Cohomologie des faisceaux}
\item \hyperref[sites-cohomology-section-phantom]{Cohomologie sur les sites}
\item \hyperref[dga-section-phantom]{Algèbre différentielle graduée}
\item \hyperref[dpa-section-phantom]{Algèbre à puissances divisées}
\item \hyperref[sdga-section-phantom]{Faisceaux différentiels gradués}
\item \hyperref[hypercovering-section-phantom]{Hyperrecouvrements}
\end{enumerate}
Schémas
\begin{enumerate}
\setcounter{enumi}{25}
\item \hyperref[schemes-section-phantom]{Schémas}
\item \hyperref[constructions-section-phantom]{Constructions de schémas}
\item \hyperref[properties-section-phantom]{Propriétés des schémas}
\item \hyperref[morphisms-section-phantom]{Morphismes de schémas}
\item \hyperref[coherent-section-phantom]{Cohomologie des schémas}
\item \hyperref[divisors-section-phantom]{Diviseurs}
\item \hyperref[limits-section-phantom]{Limites de schémas}
\item \hyperref[varieties-section-phantom]{Variétés}
\item \hyperref[topologies-section-phantom]{Topologies sur les schémas}
\item \hyperref[descent-section-phantom]{Descente}
\item \hyperref[perfect-section-phantom]{Catégories dérivées de schémas}
\item \hyperref[more-morphisms-section-phantom]{Compléments sur les morphismes}
\item \hyperref[flat-section-phantom]{Compléments sur la platitude}
\item \hyperref[groupoids-section-phantom]{Schémas en groupoïdes}
\item \hyperref[more-groupoids-section-phantom]{Compléments sur les schémas en groupoïdes}
\item \hyperref[etale-section-phantom]{Morphismes étales de schémas}
\end{enumerate}
Thèmes de théorie des schémas
\begin{enumerate}
\setcounter{enumi}{41}
\item \hyperref[chow-section-phantom]{Homologie de Chow}
\item \hyperref[intersection-section-phantom]{Théorie de l'intersection}
\item \hyperref[pic-section-phantom]{Schémas de Picard des courbes}
\item \hyperref[weil-section-phantom]{Théories cohomologiques de Weil}
\item \hyperref[adequate-section-phantom]{Modules adéquats}
\item \hyperref[dualizing-section-phantom]{Complexes dualisants}
\item \hyperref[duality-section-phantom]{Dualité pour les schémas}
\item \hyperref[discriminant-section-phantom]{Discriminants et différentes}
\item \hyperref[derham-section-phantom]{Cohomologie de de Rham}
\item \hyperref[local-cohomology-section-phantom]{Cohomologie locale}
\item \hyperref[algebraization-section-phantom]{Géométrie algébrique et formelle}
\item \hyperref[curves-section-phantom]{Courbes algébriques}
\item \hyperref[resolve-section-phantom]{Résolution des surfaces}
\item \hyperref[models-section-phantom]{Réduction semi-stable}
\item \hyperref[functors-section-phantom]{Foncteurs et morphismes}
\item \hyperref[equiv-section-phantom]{Catégories dérivées de variétés}
\item \hyperref[pione-section-phantom]{Groupes fondamentaux de schémas}
\item \hyperref[etale-cohomology-section-phantom]{Cohomologie étale}
\item \hyperref[crystalline-section-phantom]{Cohomologie cristalline}
\item \hyperref[proetale-section-phantom]{Cohomologie pro-étale}
\item \hyperref[relative-cycles-section-phantom]{Cycles relatifs}
\item \hyperref[more-etale-section-phantom]{Compléments de cohomologie étale}
\item \hyperref[trace-section-phantom]{La formule des traces}
\end{enumerate}
Espaces algébriques
\begin{enumerate}
\setcounter{enumi}{64}
\item \hyperref[spaces-section-phantom]{Espaces algébriques}
\item \hyperref[spaces-properties-section-phantom]{Propriétés des espaces algébriques}
\item \hyperref[spaces-morphisms-section-phantom]{Morphismes d'espaces algébriques}
\item \hyperref[decent-spaces-section-phantom]{Espaces algébriques décents}
\item \hyperref[spaces-cohomology-section-phantom]{Cohomologie des espaces algébriques}
\item \hyperref[spaces-limits-section-phantom]{Limites d'espaces algébriques}
\item \hyperref[spaces-divisors-section-phantom]{Diviseurs sur les espaces algébriques}
\item \hyperref[spaces-over-fields-section-phantom]{Espaces algébriques sur des corps}
\item \hyperref[spaces-topologies-section-phantom]{Topologies sur les espaces algébriques}
\item \hyperref[spaces-descent-section-phantom]{Descente et espaces algébriques}
\item \hyperref[spaces-perfect-section-phantom]{Catégories dérivées d'espaces}
\item \hyperref[spaces-more-morphisms-section-phantom]{Compléments sur les morphismes d'espaces}
\item \hyperref[spaces-flat-section-phantom]{Platitude sur les espaces algébriques}
\item \hyperref[spaces-groupoids-section-phantom]{Groupoïdes dans les espaces algébriques}
\item \hyperref[spaces-more-groupoids-section-phantom]{Compléments sur les groupoïdes dans les espaces}
\item \hyperref[bootstrap-section-phantom]{Amorçage}
\item \hyperref[spaces-pushouts-section-phantom]{Sommes amalgamées d'espaces algébriques}
\end{enumerate}
Thèmes de géométrie
\begin{enumerate}
\setcounter{enumi}{81}
\item \hyperref[spaces-chow-section-phantom]{Groupes de Chow des espaces}
\item \hyperref[groupoids-quotients-section-phantom]{Quotients de groupoïdes}
\item \hyperref[spaces-more-cohomology-section-phantom]{Compléments sur la cohomologie des espaces}
\item \hyperref[spaces-simplicial-section-phantom]{Espaces simpliciaux}
\item \hyperref[spaces-duality-section-phantom]{Dualité pour les espaces}
\item \hyperref[formal-spaces-section-phantom]{Espaces algébriques formels}
\item \hyperref[restricted-section-phantom]{Algébrisation des espaces formels}
\item \hyperref[spaces-resolve-section-phantom]{Résolution des surfaces revisitée}
\end{enumerate}
Théorie des déformations
\begin{enumerate}
\setcounter{enumi}{89}
\item \hyperref[formal-defos-section-phantom]{Théorie formelle des déformations}
\item \hyperref[defos-section-phantom]{Théorie des déformations}
\item \hyperref[cotangent-section-phantom]{Le complexe cotangent}
\item \hyperref[examples-defos-section-phantom]{Problèmes de déformation}
\end{enumerate}
Champs algébriques
\begin{enumerate}
\setcounter{enumi}{93}
\item \hyperref[algebraic-section-phantom]{Champs algébriques}
\item \hyperref[examples-stacks-section-phantom]{Exemples de champs}
\item \hyperref[stacks-sheaves-section-phantom]{Faisceaux sur les champs algébriques}
\item \hyperref[criteria-section-phantom]{Critères de représentabilité}
\item \hyperref[artin-section-phantom]{Axiomes d'Artin}
\item \hyperref[quot-section-phantom]{Espaces de Quot et de Hilbert}
\item \hyperref[stacks-properties-section-phantom]{Propriétés des champs algébriques}
\item \hyperref[stacks-morphisms-section-phantom]{Morphismes de champs algébriques}
\item \hyperref[stacks-limits-section-phantom]{Limites de champs algébriques}
\item \hyperref[stacks-cohomology-section-phantom]{Cohomologie des champs algébriques}
\item \hyperref[stacks-perfect-section-phantom]{Catégories dérivées de champs}
\item \hyperref[stacks-introduction-section-phantom]{Introduction aux champs algébriques}
\item \hyperref[stacks-more-morphisms-section-phantom]{Compléments sur les morphismes de champs}
\item \hyperref[stacks-geometry-section-phantom]{La géométrie des champs}
\end{enumerate}
Thèmes de théorie des espaces de modules
\begin{enumerate}
\setcounter{enumi}{107}
\item \hyperref[moduli-section-phantom]{Champs de modules}
\item \hyperref[moduli-curves-section-phantom]{Espaces de modules de courbes}
\end{enumerate}
Divers
\begin{enumerate}
\setcounter{enumi}{109}
\item \hyperref[examples-section-phantom]{Exemples}
\item \hyperref[exercises-section-phantom]{Exercices}
\item \hyperref[guide-section-phantom]{Guide bibliographique}
\item \hyperref[desirables-section-phantom]{Desiderata}
\item \hyperref[coding-section-phantom]{Style de programmation}
\item \hyperref[obsolete-section-phantom]{Obsolète}
\item \hyperref[fdl-section-phantom]{Licence de documentation libre GNU}
\item \hyperref[index-section-phantom]{Index généré automatiquement}
\end{enumerate}
\end{multicols}
 

\bibliography{my} 
\bibliographystyle{amsalpha} 

\end{document} 
